% FST-TU_Turbulence_Skeleton_v1_de.tex
% Skeleton draft -- Anomale Dissipation und Kolmogorov-Skalierung
% Author: Lukas Geiger, 2026

\documentclass[12pt,a4paper]{article}
\usepackage[utf8]{inputenc}
\usepackage[ngerman]{babel}
\usepackage[T1]{fontenc}
\usepackage{lmodern}
\usepackage{microtype}
\usepackage{amsmath,amssymb,amsthm,mathrsfs}
\usepackage{mathtools}
\usepackage{enumitem}
\usepackage{array}
\usepackage{tabularx}
\usepackage{xcolor}
\usepackage{geometry}
\geometry{a4paper, margin=2.5cm}
\setlength{\emergencystretch}{2em}
\usepackage{tikz}
\usetikzlibrary{arrows.meta,positioning,calc}
\usepackage{hyperref}
\hypersetup{
  colorlinks=true,
  linkcolor=blue!70!black,
  citecolor=green!50!black,
  urlcolor=blue!70!black,
  pdftitle={Anomale Dissipation und die turbulente Kaskade: Eine variationelle Freie-Energie-Herleitung der Kolmogorov-Skalierung},
  pdfauthor={Lukas Geiger},
  pdfsubject={Preprint über Turbulenz, anomale Dissipation und Kolmogorov-Skalierung},
  pdfkeywords={Turbulenz, anomale Dissipation, Kolmogorov-Skalierung, Variationsprinzip, freie Energie},
  pdflang={de-DE}
}
\makeatletter
\renewcommand*{\HyperDestNameFilter}[1]{de.#1}
\makeatother

\newtheoremstyle{breakdef}
  {8pt}{8pt}
  {\normalfont}{}
  {\bfseries}{.}
  {\newline}{}
\newtheoremstyle{breakplain}
  {8pt}{8pt}
  {\itshape}{}
  {\bfseries}{.}
  {\newline}{}
\newcolumntype{Y}{>{\raggedright\arraybackslash}X}

\theoremstyle{breakdef}
\newtheorem{definition}{Definition}[section]
\theoremstyle{breakplain}
\newtheorem{theorem}[definition]{Satz}
\newtheorem{lemma}[definition]{Lemma}
\newtheorem{proposition}[definition]{Proposition}
\newtheorem{corollary}[definition]{Korollar}
\newtheorem{conjecture}[definition]{Vermutung}
\newtheorem{axiom}[definition]{Axiom}
\theoremstyle{breakdef}
\newtheorem{remark}[definition]{Bemerkung}

% ---------------------------------------------------------------------------
\title{Anomale Dissipation und die turbulente Kaskade:\\
Eine variationelle Freie-Energie-Herleitung der Kolmogorov-Skalierung}

\author{Lukas Geiger\\
\textit{Unabhängiger Forscher, Bernau im Schwarzwald, Deutschland}\\
ORCID: \href{https://orcid.org/0009-0005-7296-1534}{0009-0005-7296-1534}%
\thanks{KI-Werkzeuge (Claude von Anthropic) wurden für mathematische Formalisierung, Literaturverifikation und redaktionelle Verfeinerung eingesetzt. Alle mathematischen Inhalte wurden vom Autor entwickelt, verifiziert und freigegeben.}}

\date{März 2026}

\begin{document}
\maketitle

% ===================================================================
\begin{abstract}
Wir schlagen ein Variationsprinzip vor, das das Kolmogorov-Energiespektrum
$k^{-5/3}$ als einzigen Minimierer eines skalenaufgelösten
Freie-Energie-Funktionals liefert. Ausgehend von einer
Littlewood--Paley-Zerlegung des Geschwindigkeitsfelds definieren wir eine
relative Entropie (Kullback--Leibler-Divergenz), die die Abweichung der
spektralen Energieverteilung von einem Referenzgleichgewicht misst. Wir
führen eine breite Klasse \emph{zulässiger} Energiefluss-Operatoren ein
-- umfassend alle lokalen, dimensional konsistenten, monotonen Closures --
und minimieren die freie Energie simultan über das Energieprofil und die
Closure-Familie, unter der Nebenbedingung eines konstanten
Skala-zu-Skala-Energieflusses~$\varepsilon$. Die K41-Verteilung ist der
einzige globale Minimierer dieses simultanen Problems
(Satz~\ref{thm:K41}). Unter einer konditionalen
Entropie-Kompatibilitätshypothese über die nichtlineare Kaskade (und
einer Nicht-Entartungsbedingung an das Forcing) leiten wir eine untere
Schranke für die Dissipationsrate im Grenzwert verschwindender Viskosität
her (Satz~\ref{thm:anomalous}), was einen neuen Zugang zur anomalen
Dissipation eröffnet, der das Onsager-Vermutungsprogramm ergänzt.
Wir führen eine Hierarchie progressiv schwächerer hinreichender
Bedingungen ein -- von punktweiser Entropie-Kompatibilität über eine
skalengefensterte Variante bis zu einer dualen, falsifizierbaren
\emph{Downhill Flux Condition} (DFC), die anhand von DNS-Daten geprüft
werden kann -- und beweisen, dass diese duale DFC die benötigte
fensterbasierte Entropiekompatibilität impliziert. Die untere
Dissipationsschranke nutzt zusätzlich die ausdrücklich angenommene uniforme
Energieeinkopplung.
Intermittenz-Korrekturen vom Kolmogorov--Obukhov-Typ (K62) ergeben sich
natürlich als Fluktuationen um das variationelle Minimum über die
Hesse-Matrix des Funktionals.
\end{abstract}

\clearpage

% ===================================================================
\section{Einleitung}\label{sec:intro}

Turbulenz gehört zu den herausragenden offenen Problemen der klassischen
Physik und angewandten Mathematik. Trotz mehr als eines Jahrhunderts an
Forschung konnte bisher keine Herleitung aus ersten Prinzipien für die
statistischen Gesetze der voll entwickelten Turbulenz erzielt werden.

\subsection{Die Kolmogorov-Theorie}

In seinen bahnbrechenden Arbeiten von 1941 postulierte
Kolmogorov~\cite{K41a,K41b}, dass im Inertialbereich voll entwickelter,
homogener, isotroper Turbulenz das Energiespektrum die universelle Form
\begin{equation}\label{eq:K41}
  E(k) \;=\; C_K \,\varepsilon^{2/3}\, k^{-5/3}
\end{equation}
annimmt, wobei $\varepsilon$ die mittlere Dissipationsrate pro
Masseneinheit und $C_K \approx 1{,}5$ die Kolmogorov-Konstante bezeichnet.
Diese Vorhersage, abgeleitet aus Dimensionsanalyse und der Hypothese
lokaler Isotropie, wurde experimentell und numerisch über einen weiten
Bereich von Reynolds-Zahlen bestätigt~\cite{Frisch95,Pope00}.

\subsection{Die Onsager-Vermutung und anomale Dissipation}

Onsager~\cite{Onsager49} vermutete 1949, dass schwache Lösungen der
inkompressiblen Euler-Gleichungen mit H\"older-Regularität unterhalb
von $1/3$ kinetische Energie ohne Viskosität dissipieren können --
sogenannte \emph{anomale Dissipation}. Umgekehrt erhalten Lösungen mit
H\"older-Exponent oberhalb von $1/3$ die Energie. Die ,,positive''
Richtung wurde von Constantin, E und Titi~\cite{CET94} bewiesen; die
,,negative'' Richtung -- die Konstruktion dissipativer schwacher
Lösungen -- wurde von Isett~\cite{Isett18} mittels konvexer Integration
vervollständigt. Die Energiebilanzanalyse von Duchon und
Robert~\cite{DR00} liefert einen distributionellen Rahmen, der den
lokalen Energiedefekt mit der Regularität des Geschwindigkeitsfelds
verbindet.  Eyink~\cite{Eyink03} etablierte eine lokale Version des
Kolmogorov-Vierfünftelgesetzes, die den Zusammenhang zwischen
Regularität und Dissipation weiter beleuchtet.

\subsection{Unser Beitrag}

In dieser Arbeit verfolgen wir einen anderen Ansatz. Anstatt pathologische
schwache Lösungen zu konstruieren, zeigen wir, dass das
Kolmogorov-Spektrum die \emph{thermodynamisch bevorzugte} Verteilung von
Energie über die Skalen ist. Wir führen ein
Freie-Energie-Funktional~$\mathcal{F}$ auf dem Raum skalenaufgelöster
Energieverteilungen ein und beweisen:
\begin{enumerate}
  \item Das K41-Spektrum ist der \emph{einzige} Minimierer von $\mathcal{F}$
        unter der Nebenbedingung konstanten Energieflusses, wobei die
        Minimierung simultan über alle zulässigen lokalen Flussoperatoren
        und Energieprofile läuft (Satz~\ref{thm:K41}). Dies vermeidet
        Zirkularität: das variationelle Argument setzt keine bestimmte
        Closure voraus.
  \item Monotonie von $\mathcal{F}$ entlang von Navier--Stokes-Trajektorien
        liefert eine universelle untere Schranke für die Dissipationsrate
        im Grenzwert verschwindender Viskosität (Satz~\ref{thm:anomalous}).
\end{enumerate}

\subsection{Strukturelle Einordnung}

Im Rahmen des breiteren Programms der Freie-Energie-Spektraltheorie~\cite{FST-Unified2026}
zeigt der K41-Selektionsmechanismus die \textbf{Pattern-A-Stabilitätslogik}:
(a)~das Energiebudget führender Ordnung ist neutral (jedes
Potenzgesetzspektrum ist kinematisch zulässig); (b)~die
Konstant-Fluss-Nebenbedingung erster Ordnung schränkt ein, bestimmt aber
das Profil nicht eindeutig; (c)~die strikte Konvexität zweiter Ordnung
von~$\mathcal{F}$ selektiert K41 als einzigen Minimierer. Diese
dreistufige Hierarchie wird in Abschnitt~\ref{sec:resolvent-cascade}
diskutiert.
Unsere Ergebnisse lösen nicht das Navier--Stokes-Regularitätsproblem. Sie
geben eine unbedingte variationelle K41-Selektion und eine konditionale
Kaskadenkompatibilitätsroute für die Diskussion anomaler Dissipation; die
quantitative untere Schranke hängt weiterhin von der expliziten uniformen
Energieeinkopplungsannahme ab.

% ===================================================================
% ===================================================================
% Abbildung 1: Beweisketten-Übersicht
% ===================================================================
\begin{figure}[ht]
\centering
\begin{tikzpicture}[
  box/.style={draw, rounded corners, minimum width=3.2cm, minimum height=0.9cm,
              align=center, font=\small},
  arrow/.style={-{Stealth[length=5pt]}, thick},
  note/.style={font=\scriptsize\itshape, text=gray}
]
\node[box, fill=blue!10] (A1) at (0,0) {(A1) Lokalität\\(A3) Monotonie};
\node[box, fill=blue!10] (LP) at (5,0) {Littlewood--Paley-\\Zerlegung};
\node[box, fill=orange!15] (F) at (2.5,-1.8) {Freie Energie\\$F[E] = \sum_j \mathrm{KL}(E_j \| E_j^*)$};
\node[box, fill=green!15] (T1) at (-0.5,-3.6) {\textbf{Satz\,1:} K41\\eindeutiger Minimierer};
\node[box, fill=green!15] (T2) at (5.5,-3.6) {\textbf{Satz\,2:} Anomale\\Dissipation};
\node[box, fill=red!10] (DFC) at (5.5,-5.4) {DFC1$^\vee$\,+\,DFC2\,+\,DFC3\\(DNS/Projektionstest)};
\draw[arrow] (A1) -- (F) node[midway, left, note] {Def.\,3.1--3.2};
\draw[arrow] (LP) -- (F) node[midway, right, note] {Shell-Energie};
\draw[arrow] (F) -- (T1) node[midway, left, note] {Joint Min.};
\draw[arrow] (F) -- (T2) node[midway, right, note] {$F$-Monotonie};
\draw[arrow] (DFC) -- (T2) node[midway, right, note] {$\Rightarrow$ NL'};
\end{tikzpicture}
\caption{Logische Struktur des Beweises. Das Freie-Energie-Funktional $F[E]$
wird aus der Littlewood--Paley-Zerlegung und der zulässigen Flux-Familie
(Axiome A1, A3) konstruiert. Satz~1 (K41-Eindeutigkeit) folgt aus der
Joint-Minimierung; der konditionale Kaskadenteil von Satz~2 erfordert die
Downhill Flux Condition (DFC), während die Dissipationsuntergrenze zusätzlich
uniformen Energieinput nutzt.}
\label{fig:proof-chain}
\end{figure}

\section{Das skalenaufgelöste Freie-Energie-Funktional}\label{sec:functional}

\subsection{Littlewood--Paley-Zerlegung}

Sei $u \in L^2(\mathbb{T}^3)$ ein divergenzfreies Geschwindigkeitsfeld auf
dem Dreitorus. Wir verwenden die Standard-Littlewood--Paley-Zerlegung
\begin{equation}\label{eq:LP}
  u \;=\; \sum_{j \ge -1} \Delta_j u,
\end{equation}
wobei $\Delta_j$ der Frequenzlokalisierungsoperator bei der dyadischen
Schale $|\xi| \sim 2^j$ ist. Wir setzen $k_j = 2^j$ und definieren die
\emph{Shell-Energie}
\begin{equation}\label{eq:Ej}
  E_j \;:=\; \|\Delta_j u\|_{L^2}^2.
\end{equation}
Im Inertialbereich Übersetzt sich die K41-Vorhersage~\eqref{eq:K41} zu
\begin{equation}\label{eq:Ej_star}
  E_j^* \;=\; C_K\, \varepsilon^{2/3}\, k_j^{-5/3} \,\Delta k_j,
  \qquad \Delta k_j = k_j \ln 2 \quad (\text{dyadische Schalenbreite}).
\end{equation}

\subsection{Energiefluss-Nebenbedingung}

Der Skala-zu-Skala-Energiefluss durch die Wellenzahl $k_j$ ist über den
trilinearen Term in den Navier--Stokes-Gleichungen definiert:
\begin{equation}\label{eq:flux}
  \Pi_j \;=\; -\sum_{\ell \le j}
    \langle \Delta_\ell(u \cdot \nabla u),\; \Delta_\ell u \rangle.
\end{equation}
Im Inertialbereich erfüllt eine statistisch stationäre Kaskade
\begin{equation}\label{eq:const_flux}
  \Pi_j \;=\; \varepsilon \quad \text{für alle } j_{\mathrm{in}} \le j \le j_{\mathrm{dis}}.
\end{equation}

\subsection{Das Funktional}

\begin{definition}[Skalenaufgelöste freie Energie]\label{def:F}
Sei $\mathbf{E} = (E_j)_{j \in \mathcal{J}}$ eine Verteilung von
Shell-Energien über den Inertialbereich $\mathcal{J} = \{j_{\mathrm{in}}, \dots,
j_{\mathrm{dis}}\}$, mit $E_j > 0$. Sei $\mathbf{E}^* = (E_j^*)$ die
Kolmogorov-Referenzverteilung~\eqref{eq:Ej_star}. Das
\emph{skalenaufgelöste Freie-Energie-Funktional} ist
\begin{equation}\label{eq:F}
  \mathcal{F}[\mathbf{E}]
  \;:=\;
  \sum_{j \in \mathcal{J}}
  \Bigl[
    E_j \ln\!\frac{E_j}{E_j^*} \;-\; E_j \;+\; E_j^*
  \Bigr].
\end{equation}
\end{definition}

\begin{remark}
Der Summand $f(E_j) = E_j \ln(E_j/E_j^*) - E_j + E_j^*$ ist die
Kullback--Leibler-Divergenz (relative Entropie) des Paares $(E_j, E_j^*)$,
betrachtet als unnormierte Masse auf einem Einpunktraum. Es gilt
$f(E_j) \ge 0$ mit Gleichheit genau dann, wenn $E_j = E_j^*$. Daher ist
$\mathcal{F}[\mathbf{E}] \ge 0$ mit Gleichheit genau dann, wenn
$\mathbf{E} = \mathbf{E}^*$.
\end{remark}

\begin{figure}[ht]
\centering
\begin{tikzpicture}[scale=0.9]
\draw[-{Stealth}] (-0.3,0) -- (7,0) node[right] {$E_j / E_j^*$};
\draw[-{Stealth}] (0,-0.3) -- (0,4.5) node[above] {$f(E_j)$};
\draw[thick, blue!70!black, domain=0.08:4.5, smooth, samples=80]
  plot ({\x}, {\x*ln(\x)/ln(2.718) - \x + 1});
\fill[red] (1,0) circle (2.5pt);
\node[below, font=\small] at (1,-0.15) {$E_j^*$ (K41)};
\draw[dashed, gray] (1,0) -- (1,3.5);
\node[font=\small, text=blue!70!black, right] at (4.5,2.8) {$f = E_j\ln\frac{E_j}{E_j^*} - E_j + E_j^*$};
\draw[thick, red!60!black, dashed] (0.3,0.2) -- (0.3,3.8);
\draw[thick, red!60!black, dashed] (3.5,0.2) -- (3.5,3.8);
\node[font=\scriptsize, text=red!60!black] at (1.9,4.1) {fluss-zulässiger Bereich};
\draw[{Stealth}-{Stealth}, red!60!black] (0.3,3.9) -- (3.5,3.9);
\end{tikzpicture}
\caption{Die Freie-Energie-Dichte $f(E_j) = E_j \ln(E_j / E_j^*) - E_j + E_j^*$
als Funktion des Shell-Energieverhältnisses $E_j / E_j^*$. Das einzige globale
Minimum bei $E_j = E_j^*$ (K41-Spektrum) ist markiert. Die Flussnebenbedingung
schränkt den zulässigen Bereich ein, aber das Minimum liegt strikt darin, was
K41 als Joint-Minimierer bestätigt.}
\label{fig:FE-landscape}
\end{figure}

\subsection{Dimensionale Konsistenz}

Wir verifizieren, dass $\mathcal{F}$ dimensional konsistent ist. Jedes
$E_j$ hat die Dimension $[L^2\,T^{-2}]$ (Energie pro Masseneinheit,
integriert über ein Spektralband). Da $E_j^*$ dieselbe Dimension besitzt,
ist das Verhältnis $E_j/E_j^*$ dimensionslos, der Logarithmus ist
wohldefiniert, und jeder Summand in~\eqref{eq:F} hat die Dimension
$[L^2\,T^{-2}]$.

% ===================================================================
\section{Hauptergebnisse}\label{sec:results}

\subsection{Eindeutigkeit der Kolmogorov-Skalierung}

\begin{definition}[Zulässige Flussfamilie]\label{def:admissible}
Eine Familie von Flussoperatoren $\bigl(\Pi_j[\mathbf{E}]\bigr)_{j \in
\mathcal{J}}$ heißt \emph{zulässig}, wenn jedes $\Pi_j$ erfüllt:
\begin{enumerate}
  \item[\textup{(A1)}] \textbf{Lokalität.}\;
    $\Pi_j$ hängt nur von $E_{j-1},\, E_j,\, E_{j+1}$ ab.
  \item[\textup{(A2)}] \textbf{Dimensionale Konsistenz.}\;
    $[\Pi_j] = [E_j / \tau_j]$ mit der Eddy-Umschlagszeit
    $\tau_j = (k_j^2\,E_j)^{-1/2}$, sodass $\Pi_j$ die Form
    \begin{equation}\label{eq:closure_general}
      \Pi_j[\mathbf{E}]
      \;=\;
      k_j\,E_j^{3/2}\;\Phi_j\!\Bigl(
        \frac{E_{j-1}}{E_j},\;
        \frac{E_{j+1}}{E_j}
      \Bigr)
    \end{equation}
    hat, wobei $\Phi_j : (0,\infty)^2 \to (0,\infty)$ eine glatte
    dimensionslose Funktion ist mit $\Phi_j(r^*_{-},\, r^*_{+}) = \alpha$
    für eine universelle Konstante $\alpha > 0$, mit
    $r^*_{\pm} = E_{j \pm 1}^* / E_j^*$.
  \item[\textup{(A3)}] \textbf{Monotonie.}\;
    $\partial \Pi_j / \partial E_j > 0$ für alle $\mathbf{E} > 0$.
\end{enumerate}
Wir bezeichnen die Menge aller zulässigen Familien mit $\mathscr{A}$.
\end{definition}

\begin{proposition}[Dimensionale Herleitung des Flux-Präfaktors]\label{prop:A2-derived}
Axiom~\textup{(A2)} ist keine unabhängige Annahme, sondern eine Konsequenz von Lokalität~\textup{(A1)} und Dimensionsanalyse. Die Form $\Pi_j = k_j\,E_j^{3/2}\,\Phi_j(\cdot)$ ist der \emph{einzige} dimensional konsistente lokale Flussausdruck.
\end{proposition}

\begin{proof}
Nach dem Buckingham-$\Pi$-Theorem muss der Fluss $\Pi_j$ ein Produkt von Potenzen der verfügbaren dimensionsbehafteten Grössen $k_j$ und~$E_j$ sein, multipliziert mit einer dimensionslosen Funktion dimensionsloser Verhältnisse. Schreibe $\Pi_j = k_j^{a}\,E_j^{b}\,\Phi_j(\ldots)$ mit $[\Pi_j] = [L^2\,T^{-3}]$ (Energiefluss pro Masseneinheit), $[k_j] = [L^{-1}]$ und $[E_j] = [L^2\,T^{-2}]$.

Dimensionsabgleich: $[k_j^a\,E_j^b] = L^{-a}\,L^{2b}\,T^{-2b} = L^{2b-a}\,T^{-2b}$.

Gleichsetzen mit $L^2\,T^{-3}$:
\[
2b - a = 2, \qquad 2b = 3 \quad\Longrightarrow\quad b = \tfrac{3}{2},\;\; a = 1.
\]
Daher $\Pi_j = k_j\,E_j^{3/2}\,\Phi_j$, und nach~(A1) kann die dimensionslose Funktion $\Phi_j$ nur von den lokalen Verhältnissen $E_{j-1}/E_j$ und $E_{j+1}/E_j$ abhängen. Die Eddy-Umschlagszeit $\tau_j = (k_j^2\,E_j)^{-1/2}$ ist ebenso die einzige dimensional konsistente lokale Zeitskala.
\end{proof}

\begin{remark}
Die klassische Kolmogorov--Heisenberg-Closure $\Pi_j = \alpha\,k_j\,E_j^{3/2}$
entspricht dem Spezialfall $\Phi_j \equiv \alpha$. Definition~\ref{def:admissible}
erweitert die Modellklasse um alle physikalisch sinnvollen lokalen Closures:
EDQNM-Closures, Leiths Diffusionsmodell und trunkierte Triaden-Wechselwirkungen
erfüllen sämtlich (A1)--(A3). Die Familie $\mathscr{A}$ ist daher eine breite,
physikalisch motivierte Klasse und nicht eine einzelne algebraische Relation.
Proposition~\ref{prop:A2-derived} zeigt, dass (A2) keine extrinsische
Modellierungswahl ist, sondern eine notwendige Konsequenz der Dimensionsanalyse,
wodurch sich die Axiomenmenge auf zwei unabhängige Annahmen reduziert:
Lokalität~(A1) und Monotonie~(A3).
\end{remark}

\begin{theorem}[Eindeutigkeit der K41-Skalierung]\label{thm:K41}
Sei $\varepsilon > 0$ gegeben und sei $\mathscr{A}$ die Klasse der
zulässigen Flussfamilien (Definition~\ref{def:admissible}).
Betrachte das \emph{simultane} Minimierungsproblem
\begin{equation}\label{eq:min}
  \min_{\substack{\mathbf{E} > 0,\\ \Pi \in \mathscr{A}}}
  \;\mathcal{F}[\mathbf{E}]
  \qquad \text{unter} \quad
  \Pi_j[\mathbf{E}] = \varepsilon
  \;\;\text{für alle } j \in \mathcal{J}.
\end{equation}
Dann gilt:
\begin{enumerate}
  \item[\textup{(i)}]
    Für \emph{jede} zulässige Flussfamilie $\Pi \in \mathscr{A}$ besitzt
    die Nebenbedingung $\Pi_j[\mathbf{E}] = \varepsilon$ ein stationäres
    Profil $\mathbf{E}^{(\Pi)}$. Unter allen solchen stationären Profilen
    ist die K41-Verteilung
    $E_j^* = C_K\,\varepsilon^{2/3}\,k_j^{-5/3}\,\Delta k_j$
    der einzige globale Minimierer von~$\mathcal{F}$, und sie wird
    durch das Kolmogorov--Heisenberg-Mitglied $\Phi_j \equiv \alpha$ mit
    $C_K = \bigl(\alpha\,(\ln 2)^{3/2}\bigr)^{-2/3}$ realisiert.
  \item[\textup{(ii)}]
    Der Minimierer ist nicht-entartet: die Hesse-Matrix
    $D^2\mathcal{F}|_{\mathbf{E}^*}$, eingeschränkt auf den Tangentialraum
    der Nebenbedingungsmannigfaltigkeit (für jedes zulässige $\Pi$), ist
    strikt positiv definit.
\end{enumerate}
\end{theorem}

\begin{proof}[Beweis (Skizze)]
\textbf{Schritt 1: Stationäre Profile.}\;
Fixiere ein beliebiges $\Pi \in \mathscr{A}$. Nach den Axiomen (A2)--(A3)
ist die Abbildung $E_j \mapsto \Pi_j[\mathbf{E}]$ bei festen
Nachbarwerten $E_{j-1},\,E_{j+1}$ strikt wachsend mit Wertebereich
$(0,\infty)$.
Daher besitzt das System $\Pi_j[\mathbf{E}] = \varepsilon$ für jedes
$\Pi$ mindestens eine Lösung $\mathbf{E}^{(\Pi)} > 0$. (Eindeutigkeit
innerhalb einer gegebenen Closure folgt aus dem Satz über implizite
Funktionen und (A3), wird hier aber nicht benötigt.)

\textbf{Schritt 2: Untere Schranke via Konvexität.}\;
Da jeder Summand $f(E_j) = E_j\ln(E_j/E_j^*) - E_j + E_j^*$ strikt
konvex ist mit globalem Minimum $f(E_j^*) = 0$, gilt
\begin{equation}\label{eq:F_lower}
  \mathcal{F}[\mathbf{E}^{(\Pi)}]
  \;\ge\; 0
  \;=\; \mathcal{F}[\mathbf{E}^*],
\end{equation}
mit Gleichheit genau dann, wenn $E_j^{(\Pi)} = E_j^*$ für alle~$j$.

\textbf{Schritt 3: Erreichbarkeit.}\;
Wir verifizieren, dass $\mathbf{E}^*$ tatsächlich ein stationäres Profil
für ein $\Pi \in \mathscr{A}$ ist. Wähle die Kolmogorov--Heisenberg-Closure
$\Phi_j \equiv \alpha$. Dann gilt
\[
  \Pi_j[\mathbf{E}^*]
  \;=\; \alpha\,k_j\,(E_j^*)^{3/2}
  \;=\; \alpha\,k_j\,
    \bigl(C_K\,\varepsilon^{2/3}\,k_j^{-5/3}\,k_j\ln 2\bigr)^{3/2}
  \;=\; \alpha\,C_K^{3/2}\,(\ln 2)^{3/2}\,\varepsilon.
\]
Die Wahl $C_K = \bigl(\alpha\,(\ln 2)^{3/2}\bigr)^{-2/3}$ liefert
$\Pi_j[\mathbf{E}^*] = \varepsilon$. (Der geometrische Faktor $(\ln 2)^{3/2}$
aus der dyadischen Schalenbreite wird in die Kolmogorov-Konstante $C_K$
absorbiert, die dadurch durch den Closure-Parameter $\alpha$ bestimmt ist.)
Somit wird die untere Schranke~\eqref{eq:F_lower} angenommen.

\textbf{Schritt 4: Eindeutigkeit.}\;
Angenommen $\Pi' \in \mathscr{A}$ und $\mathbf{E}' \ne \mathbf{E}^*$ erfüllt
$\Pi'_j[\mathbf{E}'] = \varepsilon$. Dann liefert strikte Konvexität
$\mathcal{F}[\mathbf{E}'] > 0 = \mathcal{F}[\mathbf{E}^*]$, sodass
$\mathbf{E}'$ kein Minimierer von~\eqref{eq:min} sein kann. Daher ist
$\mathbf{E}^*$ die \emph{einzige} Lösung des simultanen Problems.

\smallskip\noindent
\textbf{Hinweis zur logischen Struktur.}\;
Die Eindeutigkeit in Schritt~4 folgt aus der strikten Positivität der KL-Divergenz abseits der Referenz, die für \emph{jede} Wahl von~$\mathbf{E}^*$ gilt. Der physikalisch nichttriviale Gehalt liegt in Schritt~3 (Erreichbarkeit): Es ist die Existenz einer zulässigen Closure, die $\Pi_j[\mathbf{E}^*] = \varepsilon$ realisiert, die K41 unter allen möglichen Referenzspektren auszeichnet. Ersetzt man $\mathbf{E}^*$ durch $k^{-3}$, so gäbe es keine zulässige Closure, die die Konstant-Fluss-Nebenbedingung am vermuteten Minimierer erfüllt (siehe Bemerkung~\ref{rem:circularity}).

\textbf{Schritt 5: Nicht-Entartung.}\;
Die Hesse-Matrix $\partial^2 \mathcal{F}/\partial E_j\,\partial E_\ell
= \delta_{j\ell}/E_j$ bei $\mathbf{E}^*$ ist diagonal mit Einträgen
$1/E_j^* > 0$. Für eine \emph{feste} Closure $\Pi \in \mathscr{A}$
erzwingt die Nebenbedingung $\Pi_j[\mathbf{E}] = \varepsilon$ genau
$|\mathcal{J}|$ Gleichungen für $|\mathcal{J}|$ Unbekannte, was
generisch isolierte Lösungen liefert; der Tangentialraum ist trivial
und Nicht-Entartung ist vacuous erfüllt. Der substantielle Inhalt
von~(ii) betrifft das \emph{simultane} Problem: Im erweiterten Raum
$(\mathbf{E}, \Pi)$ mit $\Pi$ variierend über die
(unendlich-dimensionale) Familie~$\mathscr{A}$ hat die
Nebenbedingungsmenge Tangentialrichtungen entlang derer $\Pi$ variiert.
Die uneingeschränkte Hesse-Matrix
$D^2_{\mathbf{E}}\mathcal{F}|_{\mathbf{E}^*}$ ist strikt positiv
definit (Diagonaleinträge $1/E_j^* > 0$), daher bleibt ihre
Einschränkung auf \emph{jeden} Teilraum positiv definit.

\medskip
\noindent\textbf{Schritt 5b: Interpretation zweiter Ordnung.}\;
Die Diagonaleinträge $1/E_j^* \propto k_j^{5/3}$ der Hesse-Matrix wachsen
mit der Wellenzahl, sodass kleinskalige Störungen des K41-Profils stärker
bestraft werden als großskalige. Diese skalenabhängige Steifigkeit ist
der variationelle Ursprung der Vorwärtskaskade: Abweichungen bei hohen
Wellenzahlen sind energetisch teuer und treiben das System auf der lokalen
Eddy-Umschlagszeitskala zurück zum K41-Profil. Die Einparameter-Freiheit
in der Kolmogorov--Heisenberg-Closure ($\Phi_j \equiv \alpha$) entspricht
einer einzelnen flachen Richtung in der $\Pi$-Faser des erweiterten Raums
$(\mathbf{E}, \Pi)$, die die strikte Positivität von
$D^2_{\mathbf{E}}\mathcal{F}$ nicht beeinträchtigt. Diese Struktur ist
eine Instanz des Second-Order Resolvent Dominance Patterns
aus~\cite{FST-Unified2026}.
\end{proof}

\begin{remark}[Diagonale Hesse-Schranke]\label{rem:hessian}
Die Hesse-Matrix von $\mathcal{F}$ am K41-Minimierer ist diagonal mit Einträgen $\partial^2\mathcal{F}/\partial E_j^2 = 1/E_j^*$, sodass der kleinste Eigenwert $\min_j(1/E_j^*) = 1/(C_K\,\varepsilon^{2/3}\,k_1^{-2/3}\,\Delta k_1) > 0$ beträgt. Dies liefert eine untere Schranke für die Spektrallücke der zugehörigen Dirichlet-Form (siehe Bemerkung~\ref{rem:sp-diagnostics}).
\end{remark}

\begin{remark}[Zur Rolle des Referenzspektrums]\label{rem:circularity}
Ein naheliegender Einwand ist, dass das KL-Funktional $\mathcal{F}$
\emph{konstruktionsbedingt} bei $\mathbf{E}^*$ minimiert wird, unabhängig
vom physikalischen Gehalt von~$\mathbf{E}^*$. Tatsächlich ist für jede
feste Closure $\Pi$ das unrestringierte Minimum von $\mathcal{F}$ immer
$\mathbf{E}^*$. Der nichttriviale Inhalt von Satz~\ref{thm:K41} ist
daher nicht, dass $\mathcal{F}$ sein Minimum bei $\mathbf{E}^*$ annimmt
-- das ist eine Konsequenz der KL-Konstruktion --, sondern vielmehr:
\begin{enumerate}
\item Unter allen stationären Profilen $\mathbf{E}^{(\Pi)}$ (eines pro
  Closure $\Pi \in \mathscr{A}$) erfüllt nur das K41-Profil
  $\mathbf{E}^*$ \emph{gleichzeitig} die Konstant-Fluss-Nebenbedingung
  und erreicht $\mathcal{F} = 0$. Jedes andere closure-kompatible Profil
  hat $\mathcal{F} > 0$.
\item Die Referenz $\mathbf{E}^*$ ist nicht willkürlich: Sie ist das
  \emph{einzige} Potenzgesetzprofil, für das die Konstant-Fluss-Nebenbedingung
  mit der Kolmogorov--Heisenberg-Closure konsistent ist. Ersetzt man
  $\mathbf{E}^*$ durch z.B.\ $k^{-3}$, so gäbe es keine zulässige Closure,
  die $\Pi_j = \varepsilon$ am vermuteten Minimierer erfüllt.
\end{enumerate}
Zusammengefasst: Der physikalische Input ist die Konstant-Fluss-Nebenbedingung;
das KL-Funktional liefert den Selektionsmechanismus unter deren Lösungsmenge.
\end{remark}

\begin{remark}[Selbstkonsistenz der K41-Referenz]
\label{rem:k41-selfconsistency}
Das K41-Spektrum wird nicht \emph{a priori} angenommen, sondern ergibt sich als einziges selbstkonsistentes Referenzmass: Jedes zulässige $E^*$, das (A1) dimensionale Konsistenz, (A3) Kaskadenkompatibilität und die DFC erfüllt, muss proportional zu $\varepsilon^{2/3} k^{-5/3}$ sein (Proposition~\ref{prop:A2-derived}). Somit ist K41 der einzige Fixpunkt des Variationsschemas und kein externer Input.
\end{remark}

\begin{remark}[Kaskade als Reaktionskette und Nash-Gleichgewicht]\label{rem:nash}
Die variationelle Struktur von Satz~\ref{thm:K41} erlaubt eine aufschlussreiche
spieltheoretische Interpretation. Man betrachte den Inertialbereich als
hierarchische \emph{Reaktionskette}: große Wirbel bei Skala $k_j$ übertragen
Energie an Wirbel bei Skala $k_{j+1}$, wobei jede Skala als ,,Spezies'' agiert,
deren Fitness durch ihren Energieanteil bestimmt wird. Die zugehörige
Replikatordynamik auf dem Simplex der normierten Shell-Energien
$p_j = E_j / \sum_\ell E_\ell$ hat die Form
\[
  \dot{p}_j
  \;=\;
  p_j\,\bigl(g_j(\mathbf{p}) - \bar{g}(\mathbf{p})\bigr),
\]
wobei $g_j$ die Netto-Energiegewinnrate bei Skala $j$ und $\bar{g}$ der
Populationsdurchschnitt ist. Das K41-Profil $\mathbf{p}^*$ ist ein
\emph{Nash-Gleichgewicht} dieser Dynamik: keine einzelne Skala kann ihren
Energieanteil durch einseitige Abweichung erhöhen. Die freie Energie
$\mathcal{F}$ dient als strikte Ljapunow-Funktion für die
Replikatorgleichung und gewährleistet asymptotische Stabilität von
$\mathbf{p}^*$.
Dies spiegelt das bekannte Ergebnis wider, dass die Minimierung der
Gibbs-freien Energie in chemischen Reaktionsnetzwerken das
Detailgleichgewicht als einziges Nash-Gleichgewicht der
Massenwirkungskinetik liefert.
\end{remark}

\subsection{Anomale Dissipation}

Wir formulieren zunächst die strukturelle Hypothese über den nichtlinearen
Energietransfer, die dem Argument der Freie-Energie-Monotonie zugrunde liegt.

\begin{definition}[Entropie-kompatible Kaskade]\label{def:NL}
Seien $E_j(t) = \|\Delta_j u(t)\|_{L^2}^2$ die Littlewood--Paley-Shell-Energien
und $T_j = \langle \Delta_j u,\, \Delta_j\bigl((u\cdot\nabla)u\bigr)\rangle$
der nichtlineare Transfer in Schale~$j$, sodass $\sum_j T_j = 0$.
Wir sagen, die Kaskade ist \emph{entropie-kompatibel} (bezüglich eines
Referenzspektrums $\mathbf{E}^*$), wenn für einen
Regularisierungsparameter $\delta > 0$ gilt:
\begin{equation}\label{eq:NL}
  \sum_{j}\,\ln\!\Bigl(\frac{E_j + \delta}{E_j^* + \delta}\Bigr)\,T_j
  \;\le\; 0
  \qquad \text{im distributionellen Sinne auf } (0,T).
\end{equation}
\end{definition}

\begin{remark}[Status der Annahme~\eqref{eq:NL}]
Die Bedingung~\eqref{eq:NL} ist \emph{keine} Konsequenz der
Navier--Stokes-Gleichungen allein. Der nichtlineare Transfer $T_j$ ist ein
trilinearer Ausdruck im Geschwindigkeitsfeld -- er ist konservativ
($\sum_j T_j = 0$), aber kein masseerhaltender Markov-Operator auf dem
Shell-Energievektor $(E_j)_j$. Insbesondere ist das Argument
,,Konvexität der KL-Divergenz unter masseerhaltenden Abbildungen'',
das~\eqref{eq:NL} in stochastischen Kaskadenmodellen (GOY/Sabra-Schalenmodelle,
Markov-multiplikative Kaskaden) gratis liefern würde, auf die
deterministische Navier--Stokes-Nichtlinearität nicht anwendbar. Die
Bedingung kann physikalisch motiviert werden: sie besagt, dass die
Energiekaskade Energie bevorzugt von Schalen mit hoher Überraschung
$\ln(E_j/E_j^*) \gg 0$ zu Schalen näher am Gleichgewicht transportiert,
d.h.\ die Kaskade wirkt als ,,Abwärtsfluss'' in der
Freie-Energie-Landschaft.
Hinweis: Der Regularisierungsparameter in~\eqref{eq:NL} wird mit $\delta$
(nicht $\varepsilon$) bezeichnet, um Verwechslungen mit der Dissipationsrate
zu vermeiden.
\end{remark}

Die punktweise Bedingung~\eqref{eq:NL} kann erheblich abgeschwächt werden.
Die folgende skalengefensterte Variante erfordert Entropie-Kompatibilität
nur \emph{im Mittel} über Inertialbereichs-Teilfenster und erlaubt lokale
Verletzungen, die sich über benachbarte Schalen aufheben.

\begin{definition}[Window-Entropie-kompatible Kaskade]\label{def:NLw}
Die Kaskade heißt \emph{window-entropie-kompatibel}, wenn für jedes
Inertialbereichs-Teilfenster $[j_1, j_2] \subset \mathcal{J}$ mit
$j_2 - j_1 \ge 2$ gilt:
\begin{equation}\label{eq:NLw}
  \sum_{j=j_1}^{j_2}
  \ln\!\Bigl(\frac{E_j + \delta}{E_j^* + \delta}\Bigr)\,T_j
  \;\le\;
  C_{\mathrm{bdry}}\,
  \bigl(|\varphi_{j_1}| + |\varphi_{j_2}|\bigr),
\end{equation}
wobei $\varphi_j := \ln(E_j/E_j^*)$ und $C_{\mathrm{bdry}} > 0$ eine
Konstante aus den verfügbaren Kommutator-/Projektionsschranken ist.
\end{definition}

\begin{proposition}[Window-NL genügt für die Kaskadenhypothese]
\label{prop:window-NL}
Satz~\ref{thm:anomalous} bleibt gültig, wenn die punktweise
Bedingung~\eqref{eq:NL} durch die Fensterbedingung~\eqref{eq:NLw}
ersetzt wird.
\end{proposition}

\begin{proof}[Beweisskizze]
Wende die Flussdarstellung (Lemma~\ref{lem:flux}) mit Gewichten
$\varphi_j = \ln\bigl((E_j+\delta)/(E_j^*+\delta)\bigr)$
eingeschränkt auf das Fenster $[j_1, j_2]$ an. Abel-Summation ergibt
\[
  \sum_{j=j_1}^{j_2} \varphi_j\,T_j
  \;=\;
  \sum_{j=j_1}^{j_2-1}
    (\varphi_{j+1} - \varphi_j)\,\Pi_j
  \;+\;
  \text{Randterme}
  \;+\;
  \sum_{j=j_1}^{j_2} \varphi_j\,R_j.
\]
Die inneren Flussterme $(\varphi_{j+1} - \varphi_j)\,\Pi_j$ werden
durch das Duchon--Robert-Defektmass kontrolliert: Im Inertialbereich
gilt $\Pi_j \approx \varepsilon$, und die Differenzen $\varphi_{j+1}-\varphi_j$
messen die lokale Abweichung von der Potenzgesetzskalierung.
Die Randterme tragen höchstens
$C_{\mathrm{bdry}}\,(|\varphi_{j_1}|+|\varphi_{j_2}|)$ bei, was
in die Forcing-Schranke in Schritt~2 des Beweises von
Satz~\ref{thm:anomalous} absorbiert wird. Überdeckung des
vollständigen Inertialbereichs durch überlappende Fenster der Breite
$\ge 2$ und Summation liefert die globale Freie-Energie-Monotonie~\eqref{eq:F_decay}
mit höchstens einer Konstantenfaktor-Modifikation der rechten Seite.
\end{proof}

\begin{remark}[Majorisations-Alternative]
\label{rem:majorisation}
Eine noch schwächere Formulierung ersetzt Entropie-Kompatibilität durch eine
\emph{Ordnungsbedingung}: Definiere $\mathbf{E} \succeq \mathbf{E}^*$
(Majorisierung), falls $\sum_{j \ge k} E_j \ge \sum_{j \ge k} E_j^*$ für
alle~$k$. Ist die nichtlineare Kaskade \emph{ordnungserhaltend} -- d.h.\
$\mathbf{E}(t) \succeq \mathbf{E}^*$ wird durch den Fluss aufrechterhalten --
dann folgt die Freie-Energie-Schranke aus der Schur-Konvexität von
$\mathcal{F}$ (die KL-Divergenz ist Schur-konvex auf
$\mathbb{R}_{>0}^n$). Dies vermeidet die logarithmische Gewichtung
vollständig und könnte für NS-Lösungen verifizierbar sein, deren
Energiespektrum bei allen Skalen ,,oberhalb von K41'' liegt -- eine
Bedingung, die mit DNS-Evidenz bei moderaten Reynolds-Zahlen konsistent ist.
\end{remark}

Wir identifizieren nun eine \emph{hinreichende Bedingung}, unter der die
Window-Entropie-Kompatibilität~\eqref{eq:NLw} ein \emph{Satz} ist, nicht
eine Annahme.

\begin{definition}[Downhill Flux Condition]\label{def:DFC}
Eine Leray--Hopf-schwache Lösung $u$ erfüllt die \emph{Downhill Flux
Condition} (DFC) auf dem Inertialbereichs-Teilfenster
$[j_1, j_2] \subset \mathcal{J}$, wenn:
\begin{enumerate}
  \item[\textup{(DFC1$^\vee$)}] \textbf{Duale Vorwärtskaskade.}\;
    Setze $a_j := \varphi_j-\varphi_{j+1}$ für
    $j=j_1,\ldots,j_2-1$. Es existieren Konstanten
    $\Pi_0>0$ und $C_{\mathrm{proj}}\ge0$, sodass
    \begin{equation}\label{eq:DFC1-dual}
      \sum_{j=j_1}^{j_2-1} a_j\,\Pi_j
      \;\ge\;
      \Pi_0\sum_{j=j_1}^{j_2-1} a_j
      -
      C_{\mathrm{proj}}\,
      \bigl(|\varphi_{j_1}|+|\varphi_{j_2}|\bigr).
    \end{equation}
    Die stärkere punktweise Bedingung $\Pi_j\ge\Pi_0$ für alle
    $j\in[j_1,j_2-1]$ impliziert~\eqref{eq:DFC1-dual} mit
    $C_{\mathrm{proj}}=0$, aber die duale Bedingung ist die operative
    Bedingung im Abel-Summationsargument unten.
  \item[\textup{(DFC2)}] \textbf{Spektrale Steilheit.}\;
    Das kompensierte Verhältnis $\varphi_j = \ln(E_j/E_j^*)$ ist
    monoton fallend:
    $\varphi_{j+1} \le \varphi_j$ für alle $j \in [j_1, j_2-1]$.
  \item[\textup{(DFC3)}] \textbf{Restkontrolle.}\;
    Der Kommutator-Rest $R_j$ aus der
    Flussdarstellung~\eqref{eq:flux_rep} erfüllt auf $[j_1, j_2]$:
    \begin{equation}\label{eq:DFC3-concrete}
      \sum_{j=j_1}^{j_2} |\varphi_j\,R_j|
      \;\le\;
      C_{\mathrm{rem}}\,
      \bigl(|\varphi_{j_1}| + |\varphi_{j_2}|\bigr).
    \end{equation}
    Dies ist eine quantitative Kommutator-Lokalisierungsannahme. Sie folgt
    unter geeigneten Besov-/Lokalisierungsschranken wie
    $u \in L^3_t B^{1/3}_{3,\infty}(\mathbb{T}^3)$ zusammen mit
    DFC2-Monotonie, ist aber für allgemeine Leray--Hopf-Lösungen nicht
    automatisch erfüllt.
\end{enumerate}
\end{definition}

\noindent
\textbf{Vorzeichenkonvention (durchgehend verwendet).}\;
$\Pi_j$ ist der kumulative Energiefluss durch Schale~$j$
(Gl.~\eqref{eq:flux}); $\Pi_j > 0$ bezeichnet Vorwärts-Kaskade
(groß-zu-klein-skalig). $T_j$ ist der nichtlineare Transfer \emph{in}
Schale~$j$; $\sum_j T_j = 0$. Die Flussdarstellung
(Lemma~\ref{lem:flux}) lautet
$\sum_j \varphi_j T_j = +\sum_j(\varphi_{j+1}-\varphi_j)\Pi_j +
\sum_j \varphi_j R_j$, mit \textbf{positivem Vorzeichen} vor dem
Abel-Term. Mit $a_j:=\varphi_j-\varphi_{j+1}$ liefert DFC2 $a_j\ge0$,
und der innere Abel-Term ist $-\sum_j a_j\Pi_j$. Unter DFC1$^\vee$ ist
dieser Term bis auf den expliziten Projektions-Randdefekt nicht-positiv;
die ältere punktweise Bedingung $\Pi_j\ge\Pi_0$ ist nur ein starker
hinreichender Spezialfall.

\begin{proposition}[Lokale Steigungsbedingung impliziert DFC2]\label{prop:dfc2-from-k41}
Seien $E_j = |u_j|^2$ die Schalenenergien und
$E_j^* = C_K \varepsilon^{2/3} k_j^{-5/3}\Delta k_j$ das
K41-Referenzspektrum (Definition~\ref{def:F}). Definiere
$\varphi_j := \ln(E_j / E_j^*)$. Falls auf einem Inertialfenster
\[
  \frac{E_{j+1}}{E_j}
  \;\le\;
  \frac{E_{j+1}^*}{E_j^*}
  \;=\;2^{-2/3}
  \qquad (j_1 \le j < j_2),
\]
gilt, äquivalent zu lokaler Spektraldichte-Steigung höchstens $-5/3$ nach
Umrechnung zwischen Dichte und Schalenenergien, dann gilt DFC2:
$\varphi_{j+1}\le\varphi_j$.
\end{proposition}

\begin{proof}
Die Annahme ist genau
\[
  \frac{E_{j+1}/E_{j+1}^*}{E_j/E_j^*}\le 1.
\]
Logarithmieren liefert
$\ln(E_{j+1}/E_{j+1}^*)\le \ln(E_j/E_j^*)$, also DFC2. Eine uniforme
obere Schranke $E(k)\le C E^*(k)$ würde $\varphi_j$ nur nach oben
beschränken; sie erzwingt keine Monotonie der Quotienten.
\end{proof}

\begin{remark}[A-posteriori-Konsistenz von DFC2]\label{rem:empirical-reduction}
Proposition~\ref{prop:dfc2-from-k41} liefert nur einen
\emph{A-posteriori-Konsistenzcheck}: Das variationelle K41-Profil hat die
erforderliche lokale Steigung, und auf einem Fenster steilere Spektren
erfüllen ebenfalls DFC2. Die Proposition leitet DFC2 nicht aus einer bloßen
K41-Oberhülle ab. Die logische Struktur von Satz~\ref{thm:anomalous}
benötigt DFC2 daher weiterhin als Input oder als unabhängig geprüfte
DNS-/Lokalsteigungsbedingung.
\end{remark}

\begin{remark}[Duale DFC1 und das Kolmogorov-Vier-Fünftel-Gesetz]
\label{rem:dfc1-four-fifths}
Der empirische Status von DFC1$^\vee$ (Vorwärtskaskade in der
Abel-dualen Paarung) kann substanziell aufgewertet werden, indem man ihn
mit dem Kolmogorov-Vier-Fünftel-Gesetz verbindet -- dem einzigen exakten,
nichttrivialen Resultat in voll entwickelter Turbulenz:
\[
\langle (\delta_r u)^3 \rangle \;=\; -\tfrac{4}{5}\,\varepsilon\, r \qquad (r \text{ im Inertialbereich}).
\]
Diese Identität, erstmals von Kolmogorov (1941) hergeleitet und von
Duchon und Robert~\cite{DR00} rigoros als distributionelle Identität
etabliert, operiert auf zwei verschiedenen Ebenen, die sorgfältig
getrennt werden müssen:
\begin{itemize}
\item \emph{Lokaler (distributioneller) Energiefluss:} Das
Vier-Fünftel-Gesetz ist eine Aussage über die Strukturfunktion dritter
Ordnung bei Skala $r$, gültig im distributionellen Sinne. Es beweist,
dass der \emph{mittlere} Energietransfer bei Inertialbereichsskalen
vorwärts gerichtet ist.
\item \emph{Littlewood--Paley-Flusspaarung:} Punktweise
Schalenpositivität $\Pi_j \geq \Pi_0>0$ ist strikt stärker als diese
distributionelle Aussage. Die im Beweis verwendete Größe ist stattdessen
die gewichtete Paarung $\sum_j(\varphi_j-\varphi_{j+1})\Pi_j$.
\end{itemize}
Die Auflösung besteht darin, dass das Variationsargument von
Theorem~\ref{thm:K41} \emph{nicht} punktweises DFC1 erfordert. Hinreichend
ist die duale Fenster-Ungleichung
\[
\sum_{j=j_1}^{j_2-1}(\varphi_j-\varphi_{j+1})\Pi_j
\;\ge\;
\Pi_0\sum_{j=j_1}^{j_2-1}(\varphi_j-\varphi_{j+1})
- C_{\mathrm{proj}}(|\varphi_{j_1}|+|\varphi_{j_2}|).
\]
Das Vier-Fünftel-Gesetz, Duchon--Robert und Eyink liefern den richtigen
grobskaligen Flussanker für diese Bedingung, beweisen aber nicht von
selbst die Littlewood--Paley-Projektionsabschätzung. Die verbleibende
Brücke ist eine kontrollierte Projektions-/Koarea-Abschätzung von positivem
grobskaligem Skalenfluss zu DFC1$^\vee$. DFC1$^\vee$ ist hier daher kein
unbedingter Satz, sondern die scharfe falsifizierbare Bedingung, welche
die ältere punktweise Shell-Forderung ersetzt.
\end{remark}

\begin{lemma}[Downhill Flux impliziert Window-NL]\label{lem:DFC-implies-NLw}
Erfüllt eine Leray--Hopf-Lösung \textup{(DFC1$^\vee$), (DFC2), (DFC3)}
auf einem
Teilfenster $[j_1, j_2]$ mit $j_2 - j_1 \ge 2$, so gilt
\begin{equation}\label{eq:DFC-quantitative}
  \sum_{j=j_1}^{j_2} \varphi_j\,T_j
  \;\le\;
  \underbrace{-\,\Pi_0
    \sum_{j=j_1}^{j_2-1}(\varphi_j - \varphi_{j+1})}_{\le\;
    -\,\Pi_0\,(\varphi_{j_1}-\varphi_{j_2})\;\le\; 0}
  \;+\;
  C_{\mathrm{bdry}}\,
  \bigl(|\varphi_{j_1}| + |\varphi_{j_2}|\bigr),
\end{equation}
wobei
\[
  C_{\mathrm{bdry}}
  \;\le\;
  C_{\mathrm{rem}}(B)+\sup_j|\Pi_j|+C_{\mathrm{proj}}.
\]
Insbesondere
gilt die Window-Entropie-Kompatibilität~\eqref{eq:NLw}. Ist die
spektrale Steilheit strikt ($\varphi_j - \varphi_{j+1} \ge \eta > 0$
für ein $\eta$), so erfüllt der negative Hauptterm bis auf denselben
Projektions-Randdefekt
\begin{equation}\label{eq:bulk-strict}
  I_1 \;\le\; -\,\eta\,\Pi_0\,(j_2 - j_1)
  + C_{\mathrm{proj}}\bigl(|\varphi_{j_1}|+|\varphi_{j_2}|\bigr).
\end{equation}
\end{lemma}

\begin{proof}
Wende die Flussdarstellung (Lemma~\ref{lem:flux} unten) mit Gewichten
$\varphi_j = \ln\bigl((E_j+\delta)/(E_j^*+\delta)\bigr)$
eingeschränkt auf $[j_1, j_2]$ an:
\begin{equation}\label{eq:DFC-decomp}
  \sum_{j=j_1}^{j_2} \varphi_j\,T_j
  \;=\;
  \underbrace{\sum_{j=j_1}^{j_2-1}
    (\varphi_{j+1} - \varphi_j)\,\Pi_j}_{I_1}
  \;+\;
  \underbrace{\varphi_{j_1}\,\Pi_{j_1-1}
    - \varphi_{j_2}\,\Pi_{j_2}}_{I_2}
  \;+\;
  \underbrace{\sum_{j=j_1}^{j_2} \varphi_j\,R_j}_{I_3}.
\end{equation}

\textbf{Term $I_1$ (innerer Fluss).}\;
Nach (DFC2) gilt $a_j:=\varphi_j-\varphi_{j+1}\ge0$. Daher
\[
  I_1
  \;=\; -\sum_{j=j_1}^{j_2-1} a_j\,\Pi_j.
\]
Mit DFC1$^\vee$ folgt
\[
  I_1
  \;\le\;
  -\,\Pi_0\sum_{j=j_1}^{j_2-1}a_j
  + C_{\mathrm{proj}}\bigl(|\varphi_{j_1}|+|\varphi_{j_2}|\bigr)
  \;\le\;
  -\,\Pi_0(\varphi_{j_1}-\varphi_{j_2})
  + C_{\mathrm{proj}}\bigl(|\varphi_{j_1}|+|\varphi_{j_2}|\bigr).
\]
Die Gleichung $\sum a_j=\varphi_{j_1}-\varphi_{j_2}$ ist dieselbe
Teleskopierung wie in der punktweisen Version.
Ist die Steilheit strikt ($\varphi_j - \varphi_{j+1} \ge \eta$),
dann
$I_1 \le -\eta\,\Pi_0\,(j_2 - j_1)
+ C_{\mathrm{proj}}(|\varphi_{j_1}|+|\varphi_{j_2}|)$.

\textbf{Term $I_2$ (Rand).}\;
Aus~\eqref{eq:DFC-decomp} folgt $I_2 = \varphi_{j_1}\,\Pi_{j_1-1}
- \varphi_{j_2}\,\Pi_{j_2}$. Da $\Pi_{j_1-1}$ und $\Pi_{j_2}$
durch $\Pi_{\max} := \sup_{j \in [j_1-1,\,j_2]}|\Pi_j|$ beschränkt sind:
\begin{equation}\label{eq:I2-bound}
  |I_2|
  \;\le\;
  \Pi_{\max}\,\bigl(|\varphi_{j_1}| + |\varphi_{j_2}|\bigr).
\end{equation}
Beachte: $\Pi_{\max}$ wird durch die Energieeinkopplung kontrolliert. In
der statistisch stationären Kaskade gilt $\Pi_j \approx \varepsilon$
gleichmässig, also $\Pi_{\max} \le C\,\varepsilon$.

\textbf{Term $I_3$ (Kommutator-Rest).}\;
Direkt nach (DFC3)~\eqref{eq:DFC3-concrete}:
\begin{equation}\label{eq:I3-bound}
  |I_3|
  \;=\;
  \biggl|\sum_{j=j_1}^{j_2} \varphi_j\,R_j\biggr|
  \;\le\;
  \sum_{j=j_1}^{j_2} |\varphi_j\,R_j|
  \;\le\;
  C_{\mathrm{rem}}\,\bigl(|\varphi_{j_1}| + |\varphi_{j_2}|\bigr).
\end{equation}
Keine weiteren Abschätzungen sind nötig: (DFC3) ist genau als die
hier verwendete Ungleichung formuliert.

\medskip\noindent
\textbf{Zusammenfassung.}\;
\[
\sum_{j=j_1}^{j_2} \varphi_j\,T_j
= I_1 + I_2 + I_3
\le
\underbrace{-\Pi_0\,(\varphi_{j_1}-\varphi_{j_2})}_{\le\,0}
+ \underbrace{(\Pi_{\max}+C_{\mathrm{rem}}+C_{\mathrm{proj}})}_
  {=:\,C_{\mathrm{bdry}}}
  (|\varphi_{j_1}|+|\varphi_{j_2}|).
\]
Dies ist genau~\eqref{eq:NLw} mit
$C_{\mathrm{bdry}} = \Pi_{\max} + C_{\mathrm{rem}} + C_{\mathrm{proj}}$.
\end{proof}

\begin{remark}[Randfälle und Summationskonventionen]
\label{rem:edge-cases}
\mbox{}
\begin{itemize}
  \item \textbf{Fensterbreite $j_2 - j_1 \ge 2$:}\;
    Die Bedingung stellt sicher, dass die innere Summe $I_1$ mindestens
    zwei Terme ($j = j_1$ und $j = j_1+1$) enthält, sodass die
    Teleskopierung $\sum_{j=j_1}^{j_2-1}(\varphi_j - \varphi_{j+1}) =
    \varphi_{j_1} - \varphi_{j_2}$ nichttrivial ist. Für $j_2 - j_1 = 1$
    hat die Summe $I_1$ einen einzigen Term und die Schranke degeneriert zu
    $|I_1| \le |\varphi_{j_1}-\varphi_{j_2}|\cdot\Pi_{\max}$, was in
    $I_2$ absorbiert wird.
  \item \textbf{Summationsgrenzen:}\;
    Innere Summe: $j = j_1, \ldots, j_2-1$ (also $j_2-j_1$ Terme).
    Randterm: $I_2 = \varphi_{j_1}\Pi_{j_1-1} -
    \varphi_{j_2}\Pi_{j_2}$ verwendet $\Pi_{j_1-1}$ (ein Index
    unterhalb des Fensters). Dies ist wohldefiniert, da der Fluss
    $\Pi_j$~\eqref{eq:flux} für alle~$j$ definiert ist.
  \item \textbf{Definition von $\Pi_{\max}$:}\;
    $\Pi_{\max} := \sup_{j \in [j_1-1,\, j_2]} |\Pi_j|$ (einschließlich
    des Randindex $j_1-1$). In der statistisch stationären
    Inertialbereichskaskade gilt $\Pi_j \approx \varepsilon$ gleichmässig,
    also $\Pi_{\max} \le (1+\delta)\,\varepsilon$ für ein kleines $\delta$
    abhängig von der Reynolds-Zahl.
\end{itemize}
\end{remark}

\begin{proposition}[Spektrale Steigung impliziert DFC2]\label{prop:slope-DFC2}
Die scharfe Schwelle für \textup{(DFC2)} ist ein dyadisches Verhältnis
$E_{j+1}/E_j \le 2^{-2/3}$ (d.h.\ eine Log-Log-Steigung $\le -2/3$),
was den Schalenbreitenfaktor $\Delta k_j = k_j \ln 2$ widerspiegelt.
Eine bequeme \emph{hinreichende} Bedingung, die direkt der
K41-Vorhersage entspricht, ist die \emph{Log-Log-Steigungsbedingung}
\begin{equation}\label{eq:slope-condition}
  \frac{\mathrm{d}}{\mathrm{d}\ln k}\,\ln E(k)
  \;\le\; -\frac{5}{3}
  \qquad \text{für } k \in [k_{j_1},\, k_{j_2}],
\end{equation}
oder äquivalent in dyadischer Form: $E_{j+1}/E_j \le 2^{-5/3}$ für
$j \in [j_1, j_2-1]$. Unter dieser Bedingung gilt \textup{(DFC2)}
mit strikter Steilheit $\varphi_j - \varphi_{j+1} \ge \ln 2 > 0$.
\end{proposition}

\begin{proof}
Da $E_j^* = C_K\,\varepsilon^{2/3}\,k_j^{-5/3}\,\Delta k_j$ mit
$\Delta k_j = k_j \ln 2$ und $k_j = 2^j$, gilt
$E_j^* = C_K\,\varepsilon^{2/3}\,(\ln 2)\,k_j^{-2/3}$, also
$E_j^*/E_{j+1}^* = (k_j/k_{j+1})^{-2/3} = 2^{2/3}$. Die
Steigungsbedingung $\frac{\mathrm{d}}{\mathrm{d}\ln k}\ln E(k) \le -5/3$
in dyadischer Form lautet $E_{j+1}/E_j \le 2^{-5/3}$, was
\emph{stärker} ist als die für (DFC2) benötigte Schwelle
$E_{j+1}/E_j \le 2^{-2/3}$. Insbesondere gilt
$E_{j+1}/E_{j+1}^* = (E_{j+1}/E_j)\cdot(E_j^*/E_{j+1}^*)\cdot(E_j/E_j^*)
\le 2^{-5/3}\cdot 2^{2/3}\cdot(E_j/E_j^*) = 2^{-1}\cdot(E_j/E_j^*)
< E_j/E_j^*$, also $\varphi_{j+1} < \varphi_j$.

\smallskip\noindent
\textbf{Scharfe Schwelle.}\;
Die minimale Steigungsbedingung für (DFC2) ist $E_{j+1}/E_j \le 2^{-2/3}$,
d.h.\ eine Log-Log-Steigung $\le -2/3$ (nicht $-5/3$). Die stärkere
Bedingung~\eqref{eq:slope-condition} wird formuliert, weil sie direkt
der K41-Vorhersage entspricht und die in DNS verifizierte Form ist.
\end{proof}

\begin{remark}[Status der Downhill Flux Condition]
\label{rem:DFC-status}
Die DFC ist eine \emph{physikalisch verifizierbare} Bedingung, keine
mathematische Annahme über die NS-Gleichungen:
\begin{itemize}
  \item \textbf{(DFC1$^\vee$) Duale Vorwärtskaskade:} DNS stützt
    positiven mittleren Energietransfer bei $R_\lambda \gtrsim 100$.
    Rückstreuung (lokalisierter inverser Transfer) tritt auf; daher ist
    punktweises $\Pi_j\ge\Pi_0$ als primitive Bedingung zu stark. Der
    relevante Test ist die gewichtete Abel-Paarung~\eqref{eq:DFC1-dual}
    ~\cite{Frisch95,ES06,Kaneda03,Gotoh02}.
  \item \textbf{(DFC2) Spektrale Steilheit:} Das kompensierte Spektrum
    $E(k)\,k^{5/3}/\varepsilon^{2/3}$ ist im Inertialbereich
    näherungsweise konstant ($\approx C_K$) und jenseits des
    Bottleneck \emph{fallend}. Bei $R_\lambda \gtrsim 200$ ist das
    Verhältnis $E_j/E_j^*$ in allen veröffentlichten DNS über den
    Inertialbereich monoton fallend~\cite{ES06,Frisch95}. Abweichungen
    von der Monotonie beschränken sich auf $\sim 2$ Schalen nahe dem
    Dissipations-Cutoff (Bottleneck-Effekt), die in die Randkonstante
    $C_{\mathrm{bdry}}$ absorbiert werden können.
  \item \textbf{(DFC3) Restkontrolle:} Die
    Ungleichung~\eqref{eq:DFC3-concrete} ist eine konditionale quantitative
    Kommutator-Lokalisierungsannahme. Sie folgt aus expliziten
    Kommutatorabschätzungen wie~\eqref{eq:remainder-bound} nur dann, wenn
    eine geeignete Besov-/Lokalisierungsschranke vorliegt, zusammen mit
    DFC2-Monotonie. Eine solche uniforme Leray--Hopf-Schranke folgt nicht
    allein aus der Energieungleichung.
\end{itemize}

\medskip\noindent
\textbf{Abhängigkeitstabelle.}
\medskip

\begin{center}
\footnotesize
\begin{tabularx}{\textwidth}{>{\raggedright\arraybackslash}p{2.5cm} Y >{\raggedright\arraybackslash}p{4.3cm} >{\raggedright\arraybackslash}p{1.9cm}}
\hline
\textbf{Label} & \textbf{Aussage} & \textbf{Typ} & \textbf{Quelle} \\
\hline
DFC1$^\vee$ & Duale Vorwärtskaskade & Empirisch/\allowbreak Projektions\-br\"ucke & \cite{Frisch95} \\
DFC2 & Spektrale Steilheit & Empirisch (DNS) & \cite{ES06} \\
DFC3 & Konditionale Kommutator-Lokalisierung & Annahme/konditionale Abschätzung & \cite{CET94} \\
Lem.~\ref{lem:DFC-implies-NLw} & DFC $\Rightarrow$ NL$'$ & \textbf{Satz} & Diese Arbeit \\
Prop.~\ref{prop:window-NL} & NL$'$ $\Rightarrow$ Satz~\ref{thm:anomalous} & \textbf{Satz} & Diese Arbeit \\
\hline
\end{tabularx}
\end{center}

\medskip\noindent
Die logische Kette lautet:
$\textup{DFC1}^{\vee} + \textup{DFC2} + \textup{DFC3}
\;\xRightarrow{\text{Lem.~\ref{lem:DFC-implies-NLw}}}\;
\textup{NL}'
\;\xRightarrow{\text{Prop.~\ref{prop:window-NL}}}\;
\textup{Satz~\ref{thm:anomalous}}$.
Die nicht-rigorosen Inputs sind DFC1$^\vee$ und DFC2, formuliert als
falsifizierbare Ungleichungen über den gemessenen gewichteten spektralen
Fluss~\eqref{eq:DFC1-dual} und die kompensierte Steigung
($\varphi_{j+1} \le \varphi_j$). Die Implikation
$\textup{DFC} \Rightarrow \textup{NL}' \Rightarrow
\textup{Satz~\ref{thm:anomalous}}$ ist rein analytisch, sobald die duale
Flussbedingung angenommen wird. Ein Gutachter kann DFC1$^\vee$--DFC2
anhand veröffentlichter DNS-Daten unabhängig vom mathematischen Rahmenwerk
verifizieren oder falsifizieren.
\end{remark}

\begin{remark}[Unabhängige Stützung der dualen DFC1 durch Vanishing-Viscosity-Theorie]
\label{rem:Drivas-VV}
Die duale Vorwärtskaskaden-Annahme (DFC1$^\vee$) erhält unabhängige theoretische
Stützung durch das Vanishing-Viscosity-Programm von
Drivas~\cite{Drivas2019cascade}:
Für jede Raum-Zeit $L^3$ schwache Lösung der Euler-Gleichungen, die als
starker Limes $d\ge 3$-dimensionaler Navier--Stokes-Lösungen entsteht, ist
das Duchon--Robert-Defektmass $D(u)$ nichtnegativ, und das Lagrangesche
Vorwärts/Rückwärts-Dispersionsmass stimmt mit dem Energiedefekt im Sinne
der Distributionen überein. Die resultierende Lagrangesche Zeitasymmetrie
(Teilchen dispergieren in $d\ge 3$ schneller rückwärts als vorwärts) ist
eine rigorose Charakterisierung der Vorwärtskaskadenrichtung.

De~Rosa, Drivas und Inversi~\cite{DeRosaDrivasInversi2024} verschärfen dies
weiter: Jede beschränkte Euler-Lösung, die als Null-Viskositäts-Limes von
Navier--Stokes-Lösungen entsteht, kann anomale Dissipation nicht auf einer
Menge mit Hausdorff-Dimension kleiner als~$d$ tragen. Diese Schranke ist scharf.

Diese Resultate beweisen die für DFC1$^\vee$ nötige
Littlewood--Paley-Projektionsbrücke \emph{nicht}, liefern aber eine
mathematisch rigorose Begründung für die physikalische Erwartung
$D(u) \ge 0$, die der gewichteten Vorwärtsfluss-Paarung zugrunde liegt.
Die quantitative gewichtete Ungleichung bleibt ein empirischer bzw.
projektiver Input.
\end{remark}

Die folgenden zwei Lemmata bilden das analytische Rückgrat. Das erste
isoliert die Kommutator-Restschranke als eigenständige Abschätzung; das
zweite liefert die Flusszerlegung, die das gesamte Argument trägt.

\begin{lemma}[Kommutator-Restschranke]\label{lem:remainder}
Sei $u$ eine Leray--Hopf-schwache Lösung auf $\mathbb{T}^3$ und sei $R_j$
der Littlewood--Paley-Kommutator-Rest aus der Zerlegung
$T_j = \Pi_{j-1} - \Pi_j + R_j$.
\begin{enumerate}
  \item[\textup{(i)}] $\sum_j R_j = 0$ (Erhaltung).
  \item[\textup{(ii)}] Gilt $u \in L^3_t B^{1/3}_{3,\infty}(\mathbb{T}^3)$
    mit Norm $\le B$, dann gilt für jedes Teilfenster $[j_1, j_2]$ und
    beliebige beschränkte Gewichte $(\varphi_j)_j$:
    \begin{equation}\label{eq:Rj-bound}
      \sum_{j=j_1}^{j_2} |\varphi_j|\,|R_j|
      \;\le\;
      C_0\,B^3\,\sup_{j \in [j_1,j_2]} |\varphi_j|,
    \end{equation}
    wobei $C_0 > 0$ eine universelle Konstante ist.
\end{enumerate}
\end{lemma}

\begin{proof}
Teil~(i) folgt aus $\sum_j T_j = 0$ und $\sum_j(\Pi_{j-1}-\Pi_j)=0$
(Teleskopierung). Teil~(ii) ist die Standard-Paraprodukt-/Kommutator-Abschätzung:
jedes $R_j$ sammelt Terme der Form $[\Delta_j, S_{j'}](u \otimes u)$,
und der Nebendiagonal-Abfall $\|[\Delta_j, S_{j'}]\|_{L^1 \to L^1}
\le C\,2^{-|j-j'|/3}$ kombiniert mit Summation über $j'$ ergibt
$\|R_j\|_{L^1} \le C_0\,B^3$; siehe~\cite{CET94}, Proposition~1
und~\cite{DR00}, Lemma~2.1. Die gewichtete Summe folgt aus
$\sum |{\varphi_j}|\,|{R_j}| \le \sup|\varphi_j|\,\sum|R_j|
  \le \sup|\varphi_j|\cdot C_0 B^3$. Es wird nicht behauptet, dass jede
Leray--Hopf-Lösung eine $\nu$-uniforme
$L^3_tB^{1/3}_{3,\infty}$-Schranke besitzt; wenn eine solche Schranke auf
einem Fenster vorliegt, gilt die obige Abschätzung.
\end{proof}

Das folgende Lemma, basierend auf der distributionellen Energiebilanz von
Duchon und Robert~\cite{DR00}, liefert die Flusszerlegung.

\begin{lemma}[Flussdarstellung]\label{lem:flux}
Für jede Familie von Gewichten $(\varphi_j)_j$ mit
$\sup_j |\varphi_j| < \infty$ erlaubt der nichtlineare Transfer die
Zerlegung
\begin{equation}\label{eq:flux_rep}
  \sum_j \varphi_j\,T_j
  \;=\;
  \sum_j (\varphi_{j+1} - \varphi_j)\,\Pi_j
  \;+\;
  \sum_j \varphi_j\,R_j,
\end{equation}
wobei $\Pi_j$ der Energiefluss durch Schale~$j$
(Gl.~\eqref{eq:flux}) ist und der Kommutator-Rest $R_j$ die
Littlewood--Paley-Paraprodukt-Kommutatorterme sammelt; $\sum_j R_j = 0$.
Für jedes Teilfenster $[j_1, j_2]$ erfüllt der gewichtete Rest die
\emph{explizite Schranke}
\begin{equation}\label{eq:remainder-bound}
  \sum_{j=j_1}^{j_2} |\varphi_j\,R_j|
  \;\le\;
  C_{\mathrm{rem}}\,\|u\|_{L^3_t B^{1/3}_{3,\infty}}^3
  \;\cdot\;
  \sup_{j \in [j_1,j_2]} |\varphi_j|,
\end{equation}
wobei $C_{\mathrm{rem}} > 0$ eine universelle Konstante ist (unabhängig
vom Fenster, der Lösung und dem Referenzspektrum). Dies folgt aus der
Standard-Paraprodukt-Zerlegung $R_j = \sum_{\ell} [T_\ell,
\Delta_j]$ und der Kommutator-Abschätzung $\|[T_\ell, \Delta_j]\|_{L^1}
\le C\,2^{-|j-\ell|/3}\,\|u\|_{B^{1/3}_{3,\infty}}^3$
(siehe~\cite{CET94}, Proposition~1; \cite{DR00}, Lemma~2.1).
\end{lemma}

\begin{proof}[Beweis (Skizze)]
Schreibe $T_j = \Pi_{j-1} - \Pi_j + R_j$, wobei
$\Pi_j = -\langle \Delta_{\le j} u \otimes \Delta_{\le j} u,\,
\nabla \Delta_{\le j} u \rangle$ der kumulative Fluss ist
und $R_j$ die Littlewood--Paley-Kommutatorterme sammelt.
Abel-Summation ergibt~\eqref{eq:flux_rep}. Die Kommutator-Abschätzung
folgt aus der Duchon--Robert-distributionellen Energiebilanz und
Standard-Paraprodukt-Abschätzungen in Besov-Räumen;
siehe~\cite{DR00} und~\cite{CET94}.
\end{proof}

\begin{theorem}[Konditionale Kaskadenkompatibilität und Energieinput-Schranke]\label{thm:anomalous}
Sei $(u^\nu)_{\nu > 0}$ eine Familie von Leray--Hopf-schwachen Lösungen
der inkompressiblen Navier--Stokes-Gleichungen auf $\mathbb{T}^3$:
\begin{equation}\label{eq:NSE}
  \partial_t u^\nu + (u^\nu \cdot \nabla)\,u^\nu
  + \nabla p^\nu
  = \nu\,\Delta u^\nu + f,
  \qquad \nabla \cdot u^\nu = 0,
\end{equation}
mit einem festen glatten Forcing $f$ und Anfangsdaten mit
$\|u^\nu_0\|_{L^2} \le M$. Bezeichne die zeitgemittelte Dissipationsrate
\begin{equation}\label{eq:diss}
  \varepsilon_\nu
  \;:=\;
  \frac{\nu}{T} \int_0^T \!\!\int_{\mathbb{T}^3}
    |\nabla u^\nu|^2 \,dx\,dt.
\end{equation}
Angenommen, die Kaskade erfüllt \emph{eine} der folgenden Bedingungen
(in absteigender Stärke):
\begin{enumerate}
  \item[\textup{(H1)}] Entropie-Kompatibilität~\eqref{eq:NL}
    (Definition~\ref{def:NL}), oder
  \item[\textup{(H2)}] Window-Entropie-Kompatibilität~\eqref{eq:NLw}
    (Definition~\ref{def:NLw}), oder
  \item[\textup{(H3)}] Die Downhill Flux Condition
    \textup{(DFC1$^\vee$), (DFC2), (DFC3)} aus Definition~\ref{def:DFC}.
\end{enumerate}
Ferner gelte \textup{($\nu$-uniforme Energieeinkopplung)}:
es existiert $m_0 > 0$, unabhängig von $\nu$, sodass
$T^{-1}\int_0^T \langle f, u^\nu \rangle\,dt \ge m_0$
für alle $\nu > 0$.

Dann liefern die Kaskadenannahmen die Entropiekompatibilitätsroute des
variationellen Rahmens. Außerdem existiert durch die globale Energiebilanz
und die uniforme Energieeinkopplungsannahme eine Konstante
$\varepsilon_* > 0$, die nur von $f$ und $M$ abhängt, sodass
\begin{equation}\label{eq:anomalous}
  \liminf_{\nu \to 0}\; \varepsilon_\nu
  \;\ge\; \varepsilon_*.
\end{equation}
\end{theorem}

\begin{proof}[Beweis (Skizze)]
\textbf{Schritt 1: Energiebilanz.}\;
Multiplikation von~\eqref{eq:NSE} mit $u^\nu$ und Integration ergibt die
globale Energiebilanz
\begin{equation}\label{eq:energy_balance}
  \frac{d}{dt}\frac{1}{2}\|u^\nu\|_{L^2}^2
  \;=\;
  -\nu\|\nabla u^\nu\|_{L^2}^2
  \;+\;
  \langle f, u^\nu \rangle.
\end{equation}
Im statistisch stationären Zustand entspricht die zeitgemittelte
Dissipation dem zeitgemittelten Energieeintrag:
$\varepsilon_\nu = T^{-1}\int_0^T \langle f, u^\nu \rangle\,dt$.

\textbf{Schritt 2: Regularisierte Freie-Energie-Monotonie.}\;
Für $\delta > 0$ definiere die regularisierte freie Energie
\[
  \mathcal{F}_\delta[\mathbf{E}]
  \;:=\;
  \sum_j \Bigl[
    (E_j + \delta)\,\ln\frac{E_j + \delta}{E_j^* + \delta}
    \;-\; (E_j - E_j^*)
  \Bigr].
\]
Diese ist $C^1$ in allen Argumenten und erfüllt $\mathcal{F}_\delta \ge 0$
mit Gleichheit genau dann, wenn $\mathbf{E} = \mathbf{E}^*$.
Berechnung von $\frac{d}{dt}\mathcal{F}_\delta$ entlang der
Shell-Energieentwicklung $\frac{1}{2}\frac{d}{dt}E_j = -\nu k_j^2 E_j + T_j
+ \langle \Delta_j u, \Delta_j f\rangle$ (bis auf Kommutatorkorrekturen)
ergibt drei Beiträge:
\begin{enumerate}
  \item \emph{Viskoser Term.} Er trägt
    \[
      -2\nu \sum_j k_j^2\,E_j^\nu\,
      \ln\!\bigl(\tfrac{E_j^\nu+\delta}{E_j^*+\delta}\bigr)
    \]
    bei. Einzelne Summanden $k_j^2\,E_j\,\ln(E_j/E_j^*)$ können positiv
    oder negativ sein: positiv bei $E_j > E_j^*$ und negativ bei
    $E_j < E_j^*$. Der viskose Beitrag hat daher kein definites Vorzeichen.
    Seine Rolle in~\eqref{eq:F_decay} ist stattdessen dissipativ: Er treibt
    das Spektrum gegen~$\mathbf{E}^*$, weil Schalen mit $E_j > E_j^*$
    schneller als im Gleichgewicht dissipiert werden, während Schalen mit
    $E_j < E_j^*$ langsamer dissipieren;
  \item \emph{Nichtlinearer Transfer.} Nach Annahme~\eqref{eq:NL}
    (entropie-kompatible Kaskade) gilt
    \[
      \sum_j \ln\bigl(\tfrac{E_j+\delta}{E_j^*+\delta}\bigr)\,T_j \le 0;
    \]
  \item \emph{Forcing:} beschränkt durch $C\,\|f\|_{H^s}^2$ via
    Cauchy--Schwarz.
\end{enumerate}
Zusammenfügen und Grenzübergang $\delta \downarrow 0$ via Fatou-Lemma
(Unterhalbstetigkeit von $\mathcal{F}$) ergibt
\begin{equation}\label{eq:F_decay}
  \frac{d}{dt}\,\mathcal{F}[\mathbf{E}^\nu(t)]
  \;\le\;
  -\,2\nu \sum_{j} k_j^2\,
    E_j^\nu\,\ln\!\frac{E_j^\nu}{E_j^*}
  \;+\;
  C\,\|f\|_{H^s}^2.
\end{equation}
Die Grösse $D_\mathcal{F}^\nu := \sum_j k_j^2\,
E_j^\nu\,\ln(E_j^\nu/E_j^*)$ ist die \emph{Entropie-Dissipation}.
Sie hat kein definites Vorzeichen (einzelne Terme können positiv oder
negativ sein). Jeder Summand $k_j^2\,E_j\,\ln(E_j/E_j^*)$ verschwindet
genau dann, wenn $E_j = E_j^*$ (da $\ln(x/y) = 0$ genau bei $x = y$
für $x,y > 0$ und $k_j^2\,E_j > 0$); folglich gilt
$D_\mathcal{F}^\nu = 0$ mit allen Summanden einzeln Null genau dann,
wenn $\mathbf{E}^\nu = \mathbf{E}^*$. Der dissipative Charakter des viskosen Terms
ist in der Gesamtungleichung~\eqref{eq:F_decay} kodiert: zusammen mit
dem nicht-positiven nichtlinearen Beitrag und dem beschränkten Forcing
stellt er sicher, dass $\mathcal{F}$ im Mittel abnimmt.

\textbf{Schritt 3: Untere Schranke für die Dissipation (mit ND).}\;
Unter der $\nu$-uniformen Energieeinkopplungsannahme (formuliert in
Schritt~4 unten) liefert die Energiebilanz (Schritt~1) direkt
$\varepsilon_\nu \ge m_0 > 0$. Ohne diese Annahme kann man nur
heuristisch argumentieren: Falls $\varepsilon_\nu \to 0$, bleibt für
glattes, nichttriviales $f$ und beschränktes $\|u^\nu\|_{L^2} \le C(M)$
der Forcing-Term $\langle f, u^\nu \rangle$ beschränkt, während die
Dissipation verschwindet, was mit nachhaltiger Energieeinkopplung
unvereinbar ist. Die ND-Annahme formalisiert dies.

\textbf{Schritt 4: Quantitative Schranke.}\;
Die $\mathcal{F}$-Monotonie~\eqref{eq:F_decay} kontrolliert die
Entropie-Dissipation $D_\mathcal{F}^\nu(t) := \sum_j k_j^2\,
E_j^\nu\,\ln(E_j^\nu/E_j^*)$, die den spektralen Abstand
von der Kolmogorov-Referenz~$\mathbf{E}^*$ misst.
Allerdings wird die physikalische Dissipation
$\varepsilon_\nu = \nu\sum_j k_j^2 E_j^\nu$ nicht direkt
durch~$D_\mathcal{F}^\nu$ kontrolliert, ohne eine zusätzliche
\emph{Nicht-Entartungs}-Annahme, die beide Grössen verknüpft.

\emph{Annahme ND ($\nu$-uniforme Energieeinkopplung).}\;
Es existiert $m_0 > 0$, unabhängig von $\nu$, sodass
\begin{equation}\label{eq:ND}
  \frac{1}{T}\int_0^T \langle f, u^\nu \rangle\,dt \;\ge\; m_0 > 0
  \qquad \text{für alle } \nu > 0.
\end{equation}
Unter dieser Annahme liefert die Energiebilanz (Schritt~1) unmittelbar
$\varepsilon_\nu \ge m_0$ für alle~$\nu$, also
$\liminf_{\nu\to 0}\,\varepsilon_\nu \ge m_0 > 0$.

Eine verfeinerte Schranke erhält man wie folgt. Für divergenzfreies
$f \in H^{-1}(\mathbb{T}^3)$ mit $\langle f, \cdot\rangle$ nichttrivial
auf divergenzfreien Feldern liefert die Poincar\'e-Ungleichung auf
$\mathbb{T}^3$: $\|u^\nu\|_{L^2} \le C_P\,\|\nabla u^\nu\|_{L^2}$,
sodass
$\langle f, u^\nu \rangle
\le \|f\|_{H^{-1}}\,\|\nabla u^\nu\|_{L^2}
\le \|f\|_{H^{-1}}\,(\varepsilon_\nu/\nu)^{1/2}$.
Kombination mit~\eqref{eq:ND}:
$m_0 \le \|f\|_{H^{-1}}\,(\varepsilon_\nu/\nu)^{1/2}$, was
$\varepsilon_\nu \ge \nu\,m_0^2/\|f\|_{H^{-1}}^2$ ergibt. Diese
Schranke ist nur für festes~$\nu$ nichttrivial und degeneriert für
$\nu \to 0$.
Die $\nu$-\emph{unabhängige} untere Schranke
$\varepsilon_* \ge m_0$ aus Schritt~1 bleibt die operative Abschätzung;
eine schärfere $\nu$-unabhängige Schranke der Form
$\varepsilon_* \ge c\,\|f\|_{H^{-1}}^2 / M^2$
würde eine spektrale Poincar\'e-Ungleichung erfordern, die
$D_\mathcal{F}^\nu$ mit~$\varepsilon_\nu$ verknüpft -- ein offenes
Problem (siehe die Bemerkung unten).
Die natürliche Sobolev-Norm für die Energieeinkopplung ist~$H^{-1}$
(dual zu~$\|\nabla u\|_{L^2}$); das $H^s$ mit $s > -1$
in~\eqref{eq:F_decay} stammt aus der schalenweisen Forcing-Zerlegung und
ist strikt schwächer.

% (Vollständige Beweisdetails in der erweiterten Version.)
\end{proof}

\begin{remark}[Rolle der Freie-Energie-Maschinerie -- ehrliche Einschätzung]
\label{rem:F-role}
Satz~\ref{thm:anomalous} enthält zwei logisch unabhängige Ergebnisse, die unterschieden werden sollten:
\begin{enumerate}
\item[(a)] \textbf{Quantitative Schranke (ND allein).}\;
  Die Schranke $\varepsilon_* \ge m_0$ folgt direkt aus der Energiebilanz und Annahme~ND, \emph{ohne} das Funktional~$\mathcal{F}$ oder die Hypothesen (H1)--(H3) zu benötigen. Dies ist die operative Abschätzung.
\item[(b)] \textbf{Strukturelle Attraktoreigenschaft ($\mathcal{F}$-Maschinerie).}\;
  Die $\mathcal{F}$-Monotonie~\eqref{eq:F_decay} liefert eine \emph{qualitative} Erklärung dafür, warum das K41-Profil der Attraktorzustand der Kaskade ist, und sie kontrolliert die Entropie-Dissipation $D_\mathcal{F}^\nu$, die den spektralen Abstand von K41 misst. Diese strukturelle Einsicht erfordert (H1)--(H3), liefert aber \emph{keine} schärfere quantitative Schranke als~(a) ohne eine zusätzliche spektrale Poincar\'e-Ungleichung (Vermutung~\ref{conj:spectral-poincare}).
\end{enumerate}
Die Hypothesen (H1)--(H3) sind daher logisch nur für die strukturelle Attraktor-Interpretation~(b) notwendig, nicht für die quantitative Anomale-Dissipations-Schranke~(a). Ein vollständiges quantitatives Upgrade von~(b) bleibt ein offenes Problem.
\end{remark}

\begin{remark}[Gültigkeitsbereich der Entropie-Kompatibilitätshypothese]
Satz~\ref{thm:anomalous} ist konditional auf
Annahme~\eqref{eq:NL}. Drei Szenarien, in denen diese Hypothese
bekannt oder erwartet ist:
\begin{enumerate}
\item \emph{Schalenmodelle:} In GOY- und Sabra-Schalenmodellen mit
  Markov-Transferstruktur ist der nichtlineare Term ein genuiner
  masseerhaltender Operator auf Shell-Energien, und~\eqref{eq:NL}
  folgt aus der Standard-Entropie-Dissipations-Ungleichung für relative
  Entropien unter doppelt stochastischen Abbildungen.
\item \emph{Stochastische Kaskadenmodelle:} In multiplikativen
  Kaskadenmodellen der Turbulenz (Mandelbrot, Kahane) ist die
  Energieumverteilung konstruktionsbedingt ein Markov-Kern, was~\eqref{eq:NL}
  liefert.
\item \emph{Navier--Stokes mit Abwärtsfluss:} Für NS-Lösungen, deren
  Energiefluss $\Pi_j$ eine ,,Abwärts''-Bedingung erfüllt -- Energie
  fliesst bevorzugt von überangeregten zu unterangeregten Schalen relativ
  zu~$\mathbf{E}^*$ -- liefert Lemma~\ref{lem:flux} kombiniert mit
  Abel-Summation~\eqref{eq:NL} bis auf kontrollierbare Kommutator-Reste.
  Ob diese Abwärtseigenschaft generisch für NS-Lösungen im
  Inertialbereich gilt, ist eine offene Frage, die eng mit der
  Onsager-Vermutung und der Struktur des Duchon--Robert-Defektmasses
  zusammenhängt~\cite{DR00}.
\end{enumerate}
\end{remark}

% ===================================================================
\section{Intermittenz-Korrekturen}\label{sec:intermittency}

Die K41-Theorie sagt voraus, dass die longitudinale Strukturfunktion
$p$-ter Ordnung skaliert wie
\begin{equation}\label{eq:Sp_K41}
  S_p(\ell) \;:=\;
  \bigl\langle |\delta_\ell u|^p \bigr\rangle
  \;\sim\; (\varepsilon\,\ell)^{p/3},
  \qquad
  \zeta_p^{\mathrm{K41}} = \frac{p}{3}.
\end{equation}
Experimente zeigen systematische Abweichungen: $\zeta_p < p/3$ für
$p > 3$, ein Kennzeichen von \emph{Intermittenz}.

\subsection{Fluktuationen um das variationelle Minimum}

In unserem Rahmenwerk entstehen Intermittenz-Korrekturen durch Fluktuationen
der Shell-Energien $\mathbf{E}$ um den K41-Minimierer $\mathbf{E}^*$.
Entwicklung von $\mathcal{F}$ bis zur zweiten Ordnung:
\begin{equation}\label{eq:F_expand}
  \mathcal{F}[\mathbf{E}^* + \delta\mathbf{E}]
  \;=\;
  \mathcal{F}[\mathbf{E}^*]
  \;+\;
  \frac{1}{2}\sum_{j}
    \frac{(\delta E_j)^2}{E_j^*}
  \;+\; O(|\delta\mathbf{E}|^3).
\end{equation}
Die Hesse-Matrix $H_{jj} = 1/E_j^*$ impliziert, dass Fluktuationen bei
kleinen Skalen (großes $j$, kleines $E_j^*$) stärker unterdrückt werden
als Fluktuationen bei großen Skalen. Diese Asymmetrie ist der variationelle
Ursprung der Intermittenz.

\subsection{Verbindung zum K62-Log-Normal-Modell}

Modelliert man die Fluktuationen $\delta E_j$ als gaussisch in der Variablen
$\ln(E_j/E_j^*)$, so ergeben sich die Strukturfunktions-Exponenten
\begin{equation}\label{eq:zeta_p}
  \zeta_p
  \;=\;
  \frac{p}{3}
  \;-\;
  \frac{\mu}{18}\,p(p-3),
\end{equation}
was exakt der K62-Hypothese (Kolmogorov--Obukhov) der verfeinerten
Ähnlichkeit mit Intermittenzparameter $\mu$ entspricht~\cite{K62,Frisch95}.
Der Parameter $\mu$ wird durch die Varianz von $\ln(E_j/E_j^*)$ unter dem
Gibbs-Mass $\exp(-\mathcal{F}/\Theta)$ bestimmt, wobei die effektive
Temperatur $\Theta > 0$ die Stärke der turbulenten Fluktuationen um das
K41-Gleichgewicht parametrisiert (analog zur Temperatur im kanonischen
Ensemble). Im Grenzfall $\Theta \to 0$ konzentriert sich das Gibbs-Mass
auf den K41-Minimierer und $\mu \to 0$ (keine Intermittenz); für
endliches $\Theta$ skaliert die log-normale Varianz als
$\mu \propto \Theta\,E_j^*$, was den Intermittenzparameter mit der
Intensität der Kaskadenfluktuationen verbindet.

\begin{remark}
Das She--L\'ev\^eque-Modell~\cite{SL94} schlägt die Alternative vor
\[
  \zeta_p = \frac{p}{9} + 2\Bigl(1 - \bigl(\tfrac{2}{3}\bigr)^{p/3}\Bigr),
\]
die experimentelle Daten für hohes $p$ genauer beschreibt. Die Herleitung
aus dem Freie-Energie-Rahmenwerk würde erfordern, die Fluktuationen als
Log-Poisson- statt Log-Normal-Prozess zu modellieren. Dies ist ein offenes
Problem innerhalb unseres Ansatzes.
\end{remark}

% ===================================================================
\subsection{Strukturelle Einordnung}\label{sec:resolvent-cascade}

Die simultane Minimierung aus Satz~\ref{thm:K41} zeigt eine dreistufige
Hierarchie: (i)~das Energiebudget führender Ordnung ist neutral (jedes
Potenzgesetzspektrum ist kinematisch zulässig); (ii)~die
Konstant-Fluss-Nebenbedingung erster Ordnung schränkt ein, bestimmt aber
das Profil nicht eindeutig; (iii)~die strikte Konvexität zweiter Ordnung
von~$\mathcal{F}$ selektiert K41 als einzigen Minimierer. Dieses Muster
-- bezeichnet als \emph{Second-Order Resolvent Dominance}
in~\cite{FST-Unified2026} -- hat Analoga in anderen variationellen
Selektionsproblemen und wird dort im Detail diskutiert.

% ===================================================================
\subsection{Robustheits-Leitplanken und formale Penalty-Darstellung}\label{sec:robustness}

Das strikte DFC-Rahmenwerk aus Definition~\ref{def:DFC} verlangt die duale
gewichtete untere Schranke~\eqref{eq:DFC1-dual} und exakte Monotonie des
kompensierten Spektrums. DNS-Daten zeigen kleinskalige Fluktuationen um
diese Idealisierungen, aber zwei naheliegende Erweiterungen müssen als
Leitplanken statt als neue unbedingte Theoreme behandelt werden.

\begin{proposition}[Relaxierte Flussbedingung: Geltungsbereich]
\label{prop:dfc-slack}
Sei $\mathcal{A}_{\textup{slack}}$ eine relaxierte zulässige Klasse, die
das exakte K41-Profil $\mathbf{E}^*$ enthält und das Ziel funktional
$\mathcal{F}[\mathbf{E}]$ unverändert lässt. Dann minimiert
$\mathbf{E}^*$ $\mathcal{F}$ über $\mathcal{A}_{\textup{slack}}$ weiterhin
exakt. Ein nichttriviales Slack-Theorem benötigt einen zusätzlichen
Datenfit-Term, eine perturbierte Referenz oder eine Nebenbedingung, die
$\mathbf{E}^*$ ausschließt.
\end{proposition}

\begin{proof}
Nach Konstruktion gilt $\mathcal{F}[\mathbf{E}]\ge0$ und
$\mathcal{F}[\mathbf{E}^*]=0$. Wenn $\mathbf{E}^*\in
\mathcal{A}_{\textup{slack}}$ liegt, kann kein zulässiges Profil einen
niedrigeren Zielfunktionswert haben. Das relaxierte Problem ist daher exakt
und keine echte Robustheitsabschätzung.
\end{proof}

\begin{remark}
Finite-Reynolds-Robustheit sollte daher als statistisches oder inverses
Problem formuliert werden: gemessene Spektren werden mit expliziten
Daten- und Projektionsfehlern gegen die K41-Referenz verglichen. Die
relaxierten DFC-Ungleichungen sind nützliche Diagnostik, aber für sich
allein keine Verschiebungsschranke weg von $\mathbf{E}^*$.
\end{remark}

\begin{proposition}[Formale Penalty-Darstellung]
\label{prop:minimax}
In einer endlichen Schalentrunkierung ist das Paar
$(\mathbf{E}^*,\boldsymbol{\Pi}^*)$ mit $\Pi_j^*=\varepsilon$ ein
Sattelpunkt des entkoppelten Penalty-Funktionals
\begin{equation}\label{eq:minimax}
  \mathcal{L}(\mathbf{E},\boldsymbol{\Pi})
  =
  \mathcal{F}[\mathbf{E}]
  -
  \lambda\sum_j(\Pi_j-\varepsilon)^2,
  \qquad \lambda>0,
\end{equation}
auf dem Produkt positiver Spektren und beschränkter Flussprofile. Dies ist
eine formale Penalty-Darstellung von K41 plus Konstantfluss, kein
dynamischer Zwei-Spieler-Satz.
\end{proposition}

\begin{proof}
Die Variablen $\mathbf{E}$ und $\boldsymbol{\Pi}$ sind in
\eqref{eq:minimax} entkoppelt. Der eindeutige Minimierer von
$\mathcal{F}$ ist $\mathbf{E}^*$, und der eindeutige Maximierer der
negativen quadratischen Penalty ist $\Pi_j=\varepsilon$ für jede Schale.
Die behaupteten Sattelpunkt-Ungleichungen folgen direkt.
\end{proof}

\begin{remark}
Die Minimax-Notation ist nützliche Buchhaltung für die
Konstantfluss-Penalty. Eine echte spieltheoretische Aussage bräuchte ein
gekoppeltes Auszahlungsfunktional, in dem die Flussvariablen durch die
nichtlineare Dynamik erzeugt werden und auf das gewählte Spektrum reagieren.
\end{remark}

\begin{remark}[Gewichtete $\ell^2$-Stabilität aus der Hesse-Schranke]
\label{rem:talagrand}
Das Funktional $\mathcal{F}[E]$ ist gleichmässig konvex um den
K41-Minimierer $E^*$ mit diagonaler Hesse-Matrix
\[
  D^2\mathcal{F}\big|_{E^*} = \mathrm{diag}(1/E_j^*).
\]
Eine direkte Taylor-Entwicklung, ohne Bezug auf
Optimal-Transport-Theorie, liefert die gewichtete
Stabilitätsschranke
\[
  \sum_j \frac{(E_j - E_j^*)^2}{E_j^*} \;\le\; 2\,\mathcal{F}[E]
\]
für alle zulässigen Spektren $E$ mit $E_j > 0$. Dies liefert eine von
der Reynolds-Zahl unabhängige Stabilitätsgarantie: Störungen des
K41-Profils werden in der $\ell^2(1/E_j^*)$-Norm durch den
Freie-Energie-Überschuss kontrolliert. Die Schranke ergänzt die
DFC-Slack-Abschätzung (Proposition~\ref{prop:dfc-slack}) und folgt aus
der elementaren Ungleichung
\[
  x\ln(x/y) - x + y \;\ge\; \frac{(x-y)^2}{2y}
  \qquad (x,y>0).
\]
\end{remark}

% ===================================================================
\subsection{Offene Vermutungen und intrinsische Charakterisierung}\label{sec:open-conjectures}

Das oben entwickelte variationelle Rahmenwerk legt mehrere präzise Vermutungen
nahe, die die Freie-Energie-Struktur mit klassischer Turbulenz-Phänomenologie
verbinden.

\begin{conjecture}[Spektrale Poincar\'e-Ungleichung]
\label{conj:spectral-poincare}
Für das Energiespektrum $E(k)$ einer statistisch stationären turbulenten Strömung bei Viskosität $\nu$ sind die Entropiedissipation $D_F^\nu = -dF[E]/dt$ und die physikalische Dissipation $\varepsilon_\nu = 2\nu \int k^2 E(k)\, dk$ durch eine spektrale Poincar\'e-Ungleichung verknüpft:
\[
D_F^\nu \geq \kappa_\nu \cdot \varepsilon_\nu,
\]
wobei $\kappa_\nu > 0$ die spektrale Poincar\'e-Konstante ist. Falls $\kappa_\nu$ für $\nu \to 0$ nach unten beschränkt bleibt (d.h.\ $\kappa_\nu \geq \kappa_0 > 0$), so ist der K41-Minimierer im nichtsviskosen Grenzfall dynamisch stabil: Störungen von $E^*$ werden mit einer Rate proportional zur physikalischen Dissipation abgebaut.
\end{conjecture}

\begin{remark}
Die spektrale Poincar\'e-Ungleichung verbindet die thermodynamische (entropiebasierte) und die physikalische (energiebasierte) Beschreibung der Turbulenz. Ihre Gültigkeit würde implizieren, dass das K41-Spektrum nicht bloss ein statischer Minimierer, sondern ein Attraktor der turbulenten Dynamik ist.
\end{remark}

\subsubsection*{Schwellen-Axiom und diagnostisches Framework}

Die spektrale Poincar\'e-Ungleichung lässt sich präzise als \emph{Schwellen-Axiom} formulieren -- die minimale strukturelle Bedingung, die dissipative Kaskadendynamik von anomalem Trapping trennt.

\begin{axiom}[Spektrale Poincar\'e -- SP-TU]\label{ax:sp-tu}
Sei $\mathcal{H} = \{E = (E_j)_{j \geq 1} : E_j \geq 0,\, \sum E_j < \infty\}$ der Raum der Shell-Energieverteilungen, sei $\mathcal{N} \subset \mathcal{H}$ der Unterraum der Nullmoden (Erhaltungsgrössen: Gesamtenergie $\sum E_j = \mathrm{const}$), und sei $\mu_{\mathcal{F}}$ das Gibbs-Mass $\mu_{\mathcal{F}} \propto e^{-\mathcal{F}[E]/\Theta}$. Es existiert eine Spektrallücke $\lambda > 0$, sodass
\begin{equation}\label{eq:sp-tu}
  \mathrm{Var}_{\mu_{\mathcal{F}}}(g) \leq \frac{1}{\lambda}\, \mathcal{E}_{\mathcal{F}}(g, g) \qquad \text{für alle } g \perp \mathcal{N},
\end{equation}
wobei $\mathcal{E}_{\mathcal{F}}(g,g) = \int |\nabla g|^2\, d\mu_{\mathcal{F}}$ die mit der Freie-Energie-Landschaft assoziierte Dirichlet-Form ist.
\end{axiom}

\begin{proposition}[Äquivalenzen von SP-TU]\label{prop:sp-equivalences}
Die folgenden Aussagen sind äquivalent:
\begin{enumerate}
\item[(i)] \textbf{Spektrallücke:} Axiom~\ref{ax:sp-tu} gilt mit $\lambda > 0$.
\item[(ii)] \textbf{Exponentielle Relaxation:} Für die Halbgruppe $P_t$ der Energiekaskaden-Dynamik gilt
  $\mathrm{Var}_{\mu_{\mathcal{F}}}(P_t g) \leq e^{-2\lambda t}\, \mathrm{Var}_{\mu_{\mathcal{F}}}(g)$ für alle $g \perp \mathcal{N}$.
\item[(iii)] \textbf{Rayleigh-Charakterisierung:}
  $\lambda = \inf_{g \perp \mathcal{N},\, g \neq 0} \frac{\mathcal{E}_{\mathcal{F}}(g,g)}{\mathrm{Var}_{\mu_{\mathcal{F}}}(g)}$.
\end{enumerate}
\end{proposition}

\begin{proof}[Beweisskizze]
(i)$\Leftrightarrow$(iii) ist die Standard-Variationscharakterisierung der Spektrallücke. (i)$\Rightarrow$(ii) folgt aus dem Spektralsatz für selbstadjungierte Erzeuger; (ii)$\Rightarrow$(i) durch Differentiation bei $t=0$.
\end{proof}

\begin{remark}[Diagnostische Schnittstelle]\label{rem:sp-diagnostics}
Die Rayleigh-Charakterisierung~(iii) liefert sowohl analytischen als auch numerischen Zugang zu $\lambda$:
\begin{itemize}
\item \textbf{Analytisch:} Testfunktionen der Form $g_j = \ln(E_j/E_j^*)$ liefern $\mathcal{E}_{\mathcal{F}}(g_j, g_j) / \mathrm{Var}(g_j) = 1/E_j^*$, also $\lambda \geq \min_j (1/E_j^*) = 1/(C_K \varepsilon^{2/3} k_1^{-2/3} \Delta k_1) > 0$. Dies ist die diagonale Hesse-Schranke (Bemerkung~\ref{rem:hessian}).
\item \textbf{Numerisch:} Schätzung von $\lambda$ aus dem exponentiellen Abfall der Autokorrelationen $\langle g(t)\, g(0) \rangle_{\mu_{\mathcal{F}}} \sim e^{-\lambda t}$ in Shell-Modell-Simulationen (vgl.\ Bemerkung~\ref{rem:sabra-dfc}).
\item \textbf{Blockzerlegung:} Lokale Poincar\'e-Ungleichung auf Shell-Clustern $[j_1, j_2]$ etablieren und über kontrollierte Überlappung verkleben (Zegarlinski-Block-Dynamik, analog zur Yang--Mills-Strategie im Begleitpaper \cite{Geiger2026-YM}).
\end{itemize}
\end{remark}

\begin{remark}[No-Go-Mechanismen]\label{rem:sp-nogo}
SP-TU versagt genau dann, wenn \emph{Quasi-Nullmoden} existieren: eine Folge $g_k \perp \mathcal{N}$ mit $\mathrm{Var}(g_k) = 1$ und $\mathcal{E}(g_k, g_k) \to 0$. Drei physikalische Mechanismen können solche Moden erzeugen:
\begin{enumerate}
\item \textbf{Metastabilität:} Die Kaskadendynamik wird in nahezu invarianten Regionen gefangen (z.B.\ Bottleneck-Akkumulation), was langsame Indikator-Observablen mit verschwindender Dirichlet-Form erzeugt.
\item \textbf{Mehrskalige Leckage:} Dissipation wirkt auf kleinen Skalen, während Varianz auf großen Skalen ,,parkt'', was eine defekte Körzivität erzeugt, die nicht durch ein einzelnes $\lambda$ geschlossen werden kann.
\item \textbf{Intermittenz-Bursts:} Seltene, intensive Ereignisse erzeugen schwerschwanzige Fluktuationen, die die dem Hesse-Bound zugrunde liegende Gauss-Annahme verletzen.
\end{enumerate}
Für K41-Turbulenz im Inertialbereich sollte keiner dieser Mechanismen dominieren: Die duale Vorwärtskaskade (DFC1$^\vee$) unterdrückt Trapping in der gewichteten Abel-Richtung, die Skalenlokalität (A1) verhindert Leckage, und die log-normale Intermittenz (Abschnitt~\ref{sec:intermittency}) fällt wie $1/N_{\mathrm{shells}}$ ab (Bemerkung~\ref{rem:mu-disambiguation}). Es wird daher \emph{erwartet}, dass SP-TU gilt, aber ein rigoroser Beweis steht noch aus.
\end{remark}

\begin{lemma}[SP-TU impliziert universelle Dissipationskontrolle]\label{lem:sp-dissipation}
Falls Axiom~\ref{ax:sp-tu} gilt, dann gilt für jede Störung $\delta E = E - E^*$ mit $\sum \delta E_j = 0$:
\begin{equation}
  \sum_j |\delta E_j|^2 / E_j^* \leq \frac{1}{\lambda}\, \mathcal{F}[E].
\end{equation}
Insbesondere sind die Subniveaumengen $\{E : \mathcal{F}[E] \leq C\}$ präkompakt in $\ell^1$, was die DS3-Bedingung für das Dissipative Selektionsprinzip liefert.
\end{lemma}

\begin{proof}[Beweisskizze]
Durch die Poincar\'e-Ungleichung~\eqref{eq:sp-tu} angewandt auf $g = \delta E / E^*$ wird die $L^2(\mu_{\mathcal{F}})$-Norm von $g$ durch die Dirichlet-Form kontrolliert, die in führender Ordnung gleich $\mathcal{F}[E]$ ist (Hesse-Approximation). Die Präkompaktheit folgt, da die gewichtete $\ell^2$-Schranke mit Gewichten $1/E_j^* \sim k_j^{2/3}$ den Schwanz der Energieverteilung kontrolliert.
\end{proof}

\begin{remark}[Endliches-$N$ Spektrall\"{u}cke via Shell-Modell-Markov-Struktur]\label{rem:finite-N-gap}
Für endlich-dimensionale Trunkierungen, also GOY/SABRA-Shell-Modelle
mit $N$ Schalen plus stochastischem Forcing, kann die
Shell-zu-Shell-Energietransfermatrix $T_{ij}$ zu einem stochastischen
Kern normiert werden. Ist dieser Kern irreduzibel und erfüllt er die
Dobrushin-artige Bedingung
\[
  \eta := \max_i \sum_{j \neq i} |T_{ij} - T_{ij}'| < 1,
\]
was bei hinreichend starkem Forcing relativ zur Nichtlinearität gilt,
so liefert die Standardtheorie der Markov-Ketten eine
Poincar\'e-Ungleichung mit Lücke
\[
  \Delta_{\mathrm{spektral}} \ge 1 - \eta > 0.
\]
Dies ist ein rigoroses endliches-$N$-Analogon der vermuteten spektralen
Poincar\'e-Ungleichung und stellt eine direkte strukturelle Parallele
zur Yang--Mills-Dobrushin--Zegarlinski-Maschinerie
\cite{Geiger2026-YM} her. Die Erweiterung auf den Limes $N \to \infty$
erfordert die Kontrolle von $\eta(N)$, also genau des
Turbulenz-Analogons des Yang--Mills-Kontinuumlimes.
\end{remark}

\begin{remark}[Intrinsische Charakterisierung von $E^*$ über K\'arm\'an--Howarth--Monin]
\label{rem:khm-intrinsic}
Die K\'arm\'an--Howarth--Monin-Relation
\[
  \partial_t \langle |\delta u|^2 \rangle
  + \nabla_r \cdot \langle |\delta u|^2 \delta u \rangle
  = -4\varepsilon/3 + 2\nu \Delta_r \langle |\delta u|^2 \rangle
\]
liefert einen intrinsischen Konstantfluss-Anker ohne \emph{a priori}
Annahme von K41. Im Inertialbereich $(\eta \ll r \ll L)$ verschwindet
der viskose Term, und die stationäre KHM-Relation ergibt
\[
  \langle |\delta u|^2 \delta u \rangle = -4\varepsilon r / 5,
\]
also das Vierfünftelgesetz. Dies stützt die Flusseingabe des
Variationsschemas, bestimmt aber nicht allein durch Fourier-Transformation
eindeutig das Spektrum zweiter Ordnung. Die Rückgewinnung von
$E(k)\sim \varepsilon^{2/3}k^{-5/3}$ benötigt weiterhin
Selbstähnlichkeits-, Lokalitäts- oder Schließungsannahmen. Die
KHM-Relation ist daher ein Konsistenzcheck der Referenz, kein
unabhängiger Eindeutigkeitsbeweis.
\end{remark}

\begin{conjecture}[Intermittenz-Exponenten aus Hesse-Fluktuationen]
\label{conj:intermittency-hessian}
Die anomalen Skalierungsexponenten $\zeta_p$ der Strukturfunktionen $S_p(r) = \langle |\delta u(r)|^p \rangle \sim r^{\zeta_p}$ lassen sich durch die Hesse-Matrix des Freie-Energie-Funktionals $F[E]$ am K41-Minimierer ausdrücken:
\[
\zeta_p = \frac{p}{3} - \frac{p(p-3)}{2} \sigma_H^2 + O(\sigma_H^4),
\]
wobei $\sigma_H^2 = \mathrm{tr}(\nabla^2 F[E^*]^{-1}) / N_{\mathrm{shells}}$ die normierte inverse Hesse-Spur (,,Fluktuationsstärke'') ist. Für $\sigma_H^2 \to 0$ (unendliche Reynolds-Zahl, Hesse-Dominanz) gilt $\zeta_p \to p/3$ (K41). Für endliches $\sigma_H^2 > 0$ reproduzieren die Korrekturen die She--L\'ev\^eque-Hierarchie $\zeta_p = p/9 + 2(1 - (2/3)^{p/3})$, falls $\sigma_H^2 = 2\ln(3/2)/9 \approx 0{,}090$.
\end{conjecture}

\begin{remark}
Diese Vermutung verbindet das variationelle Rahmenwerk mit der Phänomenologie der Intermittenz: Die Krümmung von $F[E]$ am Minimierer kontrolliert nicht nur die Stabilität (Proposition~\ref{prop:dfc-slack}), sondern auch die Statistik der Fluktuationen um diesen.
\end{remark}

\begin{remark}[Intermittenzparameter-Disambiguierung]\label{rem:mu-disambiguation}
Die Hesse-Spur-Vorhersage $\sigma_H^2 \sim 1/N_{\mathrm{shells}}$ quantifiziert die anomale Korrektur $\Delta\zeta_2 = \zeta_2 - 2/3$ zum Energiespektrum-Exponenten, nicht den Kolmogorov--Obukhov-Intermittenzparameter $\mu$ (der $\mu \approx 0{,}2$--$0{,}3$ erfüllt). Das She--L\'ev\^eque-Modell (1994) sagt $\zeta_p = p/9 + 2(1 - (2/3)^{p/3})$ vorher, was $\Delta\zeta_2 \approx 0{,}03$ für $N \approx 30$ aufgelöste Shells ergibt, konsistent mit unserer $1/N$-Vorhersage.
\end{remark}

% ===================================================================
\subsection*{Ergebnisklassifikation}

\begin{center}
\footnotesize
\begin{tabularx}{\textwidth}{|>{\raggedright\arraybackslash}p{3.8cm}|>{\raggedright\arraybackslash}p{4.0cm}|Y|}
\hline
\textbf{Ergebnis} & \textbf{Typ} & \textbf{Referenz} \\
\hline
\hline
\multicolumn{3}{|l|}{\textit{S\"atze (rigoros):}} \\
\hline
$\mathcal{F}[E] \geq 0$, $= 0$ gdw $E = E^*$ & Theorem & Thm.~\ref{thm:K41} \\
DFC$^\vee$ $\Rightarrow$ NL' & Lemma & Lem.~\ref{lem:DFC-implies-NLw} \\
Energieinput $\Rightarrow$ Dissipationsuntergrenze & Theorem & Thm.~\ref{thm:anomalous} \\
Lokale Steigungsbedingung impliziert DFC2 & Proposition & Prop.~\ref{prop:dfc2-from-k41} \\
\hline
\multicolumn{3}{|l|}{\textit{Empirische Hypothesen (DNS-pr\"ufbar):}} \\
\hline
DFC1$^\vee$ (duale Vorwärtskaskade) & Empirisch/\allowbreak Projektions\-br\"ucke & Def.~\ref{def:DFC}, $R_\lambda \gtrsim 100$ \\
\hline
\multicolumn{3}{|l|}{\textit{Numerisch best\"atigt:}} \\
\hline
$\mathcal{F}[E_{K41}] = 0$ (Minimierer) & Exakt & \texttt{compute\_F\_spectrum.py} \\
DFC2 (Sabra-Schalenmodell) & Schalenmodell & Bem.~\ref{rem:sabra-dfc} \\
$\mathcal{F} > 0$ f\"ur 500/500 Snapshots & Schalenmodell & Bem.~\ref{rem:sabra-dfc} \\
\hline
\multicolumn{3}{|l|}{\textit{DNS-falsifizierbare Vorhersagen:}} \\
\hline
$\mathcal{F}[E_{\mathrm{DNS}}] \to 0$ f\"ur $R_\lambda \to \infty$ & Vorhersage & Abschn.~\ref{sec:discussion} \\
$\varphi_j$ nicht-ansteigend im Inertialbereich & Vorhersage & DFC2 \\
\hline
\end{tabularx}
\end{center}

% ===================================================================
\section{Diskussion}\label{sec:discussion}

Die simultane Minimierung über alle zulässigen Closures ist keine tautologische Umformulierung von K41: Kolmogorovs ursprüngliche Herleitung stützt sich auf Dimensionsanalyse, während unser Satz zeigt, dass K41 der einzige Minimierer eines wohldefinierten Variationsproblems ist, was eine strukturelle Erklärung für die Universalität liefert.

\subsection{Bezug zur Onsager-Vermutung}

Satz~\ref{thm:anomalous} trennt zwei Fragen. Die untere
Dissipationsschranke folgt aus der globalen Energiebilanz plus der
angenommenen \(\nu\)-uniformen Energieeinkopplung. Der variationelle Beitrag
ist die konditionale Route von DFC1$^\vee$--DFC2--DFC3 zu einer
entropiekompatiblen Kaskadenstruktur und K41-Selektion. Das ist
komplementär zum Konvex-Integrations-Programm~\cite{Isett18,BV19}, aber
kein unbedingter Selektionssatz für alle schwachen Navier--Stokes-Lösungen.

\subsection{Die Rolle der Closure-Familie}

Die Umformulierung von Satz~\ref{thm:K41} über die zulässige
Flussfamilie (Definition~\ref{def:admissible}) schwächt die Abhängigkeit
von einem bestimmten Closure-Modell erheblich ab. Das K41-Spektrum wird
nicht durch eine einzelne algebraische Relation selektiert, sondern als
einziger Minimierer über \emph{alle} lokalen, dimensional konsistenten,
monotonen Flussoperatoren. Dennoch erlegt Axiom~(A2) eine strukturelle
Einschränkung auf -- den $k_j\,E_j^{3/2}$-Skalierungspräfaktor --, der
letztlich das dimensionale Argument Kolmogorovs kodiert. Die Herleitung
dieses Präfaktors aus der K\'arm\'an--Howarth--Monin-Relation oder aus
rigorosen Schranken für die Littlewood--Paley-trilineare Form bleibt ein
wichtiges offenes Problem.

\subsection{Regimeselektion und der laminar-turbulente Übergang}

Das Variationsprinzip aus Satz~\ref{thm:K41} selektiert das K41-Spektrum
\emph{innerhalb} des turbulenten Regimes, erklärt aber nicht den
laminar-turbulenten Übergang selbst. Eine natürliche Erweiterung würde
die strikte Konvexität von~$\mathcal{F}$ als Stabilitätskriterium
verwenden: unter den kinematisch zulässigen Strömungsregimen wird
dasjenige realisiert, dessen spektrale Energieverteilung ein stabiler
kritischer Punkt des Freie-Energie-Funktionals ist. Die quantitative
Untersuchung dieser Verbindung bleibt künftiger Arbeit vorbehalten.

\subsection{Limitationen}

Die wesentlichen Einschränkungen der vorliegenden Arbeit:
\begin{enumerate}
\item \textbf{Konditionalit\"at.}\par
  Satz~\ref{thm:anomalous} bleibt
  konditional auf eine Entropie-Kompatibilit\"atshypothese oder auf die
  hinreichende Hierarchie DFC1$^\vee$--DFC2--DFC3. DFC1$^\vee$ ist durch
  positiven grobskaligen Fluss motiviert, aber die Projektionsbr\"ucke vom
  Duchon--Robert/Eyink-Skalenfluss zur gewichteten Littlewood--Paley-
  Paarung wird hier nicht bewiesen.

  Die DFC3-Bedingung verlangt eine \emph{quantitative Schranke} für den
  Kommutator-Rest in Gl.~\eqref{eq:DFC3-concrete}. Eine solche Schranke
  folgt in Regularitäts-/Lokalisierungsregimen, in denen die
  Kommutatorabschätzung verfügbar ist, aber sie ist nicht automatisch aus
  der Leray--Hopf-Energieungleichung ableitbar. Ob eine nützliche
  DFC3-Schranke für allgemeine turbulente Vanishing-Viscosity-Familien
  gilt, bleibt offen.
\item \textbf{Inertialbereichs-Idealisierung.}\; Das Funktional $\mathcal{F}$
  ist auf dem Inertialbereich $\mathcal{J}$ definiert und ignoriert den
  Forcing- und Dissipationsbereich. Randeffekte bei $j_{\mathrm{in}}$ und
  $j_{\mathrm{dis}}$ (insbesondere der Bottleneck-Effekt) werden in
  Konstanten absorbiert, aber nicht quantitativ kontrolliert.
\item \textbf{Intermittenz.}\; Die K62-Verbindung
  (Abschnitt~\ref{sec:intermittency}) ist qualitativ. Die rigorose
  Herleitung von Intermittenz-Exponenten aus den Hesse-Fluktuationen ist
  offen.
\item \textbf{Eindeutigkeit der Referenz.}\; Obwohl
  Bemerkung~\ref{rem:circularity} argumentiert, dass $\mathbf{E}^*$ die
  einzige physikalisch konsistente Referenz ist, würde eine vollständig
  intrinsische Charakterisierung von $\mathbf{E}^*$ (ohne \emph{a priori}
  Bezug auf K41) das Argument stärken.
\end{enumerate}

\subsection{Numerische Evidenz und vorgeschlagener DNS-Test}

Direkte numerische Simulationen (DNS) bei Taylor-Mikroskalenreynoldszahlen
$R_\lambda \sim 200$--$1000$ erzeugen konsistent Energiespektren nahe
$k^{-5/3}$ im Inertialbereich, mit einer Kolmogorov-Konstante
$C_K \approx 1{,}5$~\cite{ES06,Frisch95,MY75}. Wir schlagen folgendes
reproduzierbares Testprotokoll vor:
\begin{enumerate}
\item Berechne die Littlewood--Paley-Shell-Energien $E_j$ aus einer
  statistisch stationären DNS bei $R_\lambda \ge 200$.
\item Werte $\mathcal{F}[\mathbf{E}]$ unter Verwendung des gemessenen
  $\mathbf{E}$ und der K41-Referenz $\mathbf{E}^*$ aus, wobei
  $\varepsilon := \nu\,\langle\|\nabla u\|_{L^2}^2\rangle_t$
  die zeitgemittelte viskose Dissipationsrate derselben DNS-Rechnung ist.
\item Verifiziere, dass $\mathcal{F}$ mit steigendem $R_\lambda$
  abnimmt und dass das kompensierte Verhältnis
  $\varphi_j = \ln(E_j/E_j^*)$ im Inertialbereich monoton fallend
  in $j$ ist (DFC2).
\item Berechne den spektralen Fluss $\Pi_j$ und verifiziere die gewichtete
  Ungleichung~\eqref{eq:DFC1-dual} für
  $a_j=\varphi_j-\varphi_{j+1}\ge0$ auf Inertialfenstern
  (DFC1$^\vee$). Punktweises $\Pi_j\ge\Pi_0>0$ ist nur ein stärkerer
  optionaler Test.
\end{enumerate}
Schritte 3--4 stellen einen direkten, falsifizierbaren Test der Downhill
Flux Condition dar, der vom mathematischen Rahmenwerk unabhängig ist.

\smallskip\noindent
\textbf{Implementierungshinweise.}\;
(a)~Standard-DNS-Codes geben das isotrope Energiespektrum $E(k)$ über
scharfes Kugelschalenbinning aus, nicht über Littlewood--Paley-gefilterte
Energien. Für glatte Spektren ist der Unterschied von der Ordnung
$O(\Delta k_j / k_j) = O(\ln 2)$ pro Schale und kann durch Vergleich
scharfen und glatten Binnings abgeschätzt werden.
(b)~Schritte 1--4 sollten auf Momentanspektren an mehreren statistisch
unabhängigen Zeitsnapshots innerhalb des stationären Regimes angewendet
werden; Ensemble-Mittelung von $\mathcal{F}$ und $\varphi_j$ liefert
dann Konfidenzintervalle für die DFC-Verifikation.

\subsection{Falsifikationsprotokoll}
\label{sec:falsification}

Die Vorhersagen dieses Rahmenwerks können anhand von Daten direkter numerischer Simulationen (DNS) wie folgt getestet werden:

\begin{enumerate}
\item \textbf{Spektraldaten:} Extrahiere das Energiespektrum $E(k)$ aus einer gut aufgelösten DNS bei Taylor-Mikroskalenreynoldszahl $Re_\lambda \geq 200$.

\item \textbf{Funktionalauswertung:} Berechne $F[E] = D_{\mathrm{KL}}(E \| E^*) + \lambda \sum_j |\Pi_j[E] - \varepsilon|$ unter Verwendung schalengemittelter Spektren mit logarithmischem Binning (Schalenverhältnis $r = 2^{1/3}$).

\item \textbf{Minimierervergleich:} Vergleiche das DNS-Spektrum $E_{\mathrm{DNS}}(k)$ mit der K41-Vorhersage $E^*(k) = C_K \varepsilon^{2/3} k^{-5/3}$ im Inertialbereich $k_f \ll k \ll k_\eta$.

\item \textbf{Falsifikationskriterium.}
  Da $\mathcal{F}[E^*] = 0$ konstruktionsbedingt gilt, kann das
  Funktional nicht durch $\mathcal{F}[E_{\mathrm{DNS}}] < 0$
  falsifiziert werden. Stattdessen ist die DFC-Route herausgefordert,
  wenn DFC1$^\vee$ oder DFC2 bei hohem $R_\lambda$ \emph{verletzt}
  werden, wo die Theorie ihr Gelten vorhersagt, oder wenn
  $\mathcal{F}[E_{\mathrm{DNS}}]$ nicht mit steigendem $R_\lambda$
  abnimmt, was darauf hindeuten würde, dass K41 nicht der asymptotische
  Attraktor ist. Punktweise Fehlschläge von $\Pi_j\ge\Pi_0$ allein
  widerlegen die duale Route nicht.

\item \textbf{Intermittenztest:} Berechne die Hesse-Eigenwerte $\lambda_j = \partial^2 F / \partial E_j^2$ am DNS-Spektrum. Negative Eigenwerte im Dissipationsbereich zeigen Instabilität des K41-Minimierers auf diesen Skalen an.
\end{enumerate}

\begin{remark}[DFC als energetische Stabilitätsklasse]\label{rem:dfc-stability}
Die Downhill-Flux-Bedingung sollte als duales/fensterbasiertes
energetisches Stabilitätsdiagnostikum geprüft werden, nicht als
punktweises Shell-Axiom. Öffentlich verfügbare DNS-Datensätze (JHTDB,
$\mathrm{Re}_\lambda \sim 200$) sind dafür geeignet, weil der Test relative
Free-Energy-Differenzen und Abel-gewichtete Flusspaarungen statt nur
asymptotischer Exponenten nutzt. Bei nicht-eindeutigen Leray--Hopf-Lösungen
wäre DFC ein falsifizierbares Selektionsdiagnostikum; dieses Paper beweist
nicht, dass DFC den physikalisch realisierten Zweig auswählt.
\end{remark}

\begin{remark}[Numerische Verifikation von Satz~\ref{thm:K41}]\label{rem:numerical-F}
Eine numerische Auswertung von $\mathcal{F}[E]$ über sechs synthetische Energiespektren bestätigt, dass das K41-Spektrum der einzige Minimierer ist: $\mathcal{F}[E_{\mathrm{K41}}] = 0$ exakt, während $\mathcal{F}[E_{\mathrm{K62}}] = 2{,}0 \times 10^{-4}$ (Intermittenzkorrektion), $\mathcal{F}[E_{k^{-2}}] = 1{,}74$ (steil), $\mathcal{F}[E_{k^{-4/3}}] = 5{,}94$ (flach) und $\mathcal{F}[E_{\mathrm{thermal}}] = 7{,}55$ (Equipartition). Ein Störungstest ($300$ zufällige Störungen $\delta E$ bei drei Rauschstärken $\|\delta E\|/\|E^*\| \in \{0{,}01,\, 0{,}1,\, 1{,}0\}$) ergibt $\mathcal{F}[E^* + \delta E] > 0$ für \emph{alle} $300$ Stichproben und bestätigt damit numerisch die strikte Konvexität. Die DFC-Bedingungen sind von K41 und K62 (mit She--L\'eveque-Intermittenz $\mu = 0{,}03$) erfüllt, werden jedoch von steileren, flacheren oder thermischen Spektren verletzt. Numerische Skripte: \texttt{compute\_F\_spectrum.py}.
\end{remark}

\begin{remark}[Sabra-Schalenmodell-Validierung der DFC]
\label{rem:sabra-dfc}
Punktweises DFC1 und DFC2 wurden am Sabra-Schalenmodell (L'vov et al., 1998) mit $N=22$ Schalen, $\nu = 10^{-7}$, unter Verwendung eines IMEX-Integrators (500 Snapshots) getestet. DFC2 (spektrale Steilheit) wird stark bestätigt: Die Inertialbereichs-Steigungen betragen $-0{,}663$, $-0{,}675$, $-0{,}661$, $-0{,}699$, $-0{,}700$ für Schalen $n=5,\ldots,9$ (K41-Vorhersage: $-2/3 \approx -0{,}667$, Abweichung $< 2\%$). Die freie Energie erfüllt $\mathcal{F}[E] > 0$ für alle 500 Snapshots (100\%), mit $\mathcal{F}[\langle E \rangle] = 0{,}198$, was das K41-Spektrum als approximativen Minimierer bestätigt. Der punktweise Fluss $\Pi_n$ ist vorzeichenalternierend, ein bekanntes Merkmal von Schalenmodellen; dies motiviert DFC1$^\vee$ als gewichtete Abel-Paarung statt punktweiser Positivität. Bestehendes Skript: \texttt{compute\_goy\_shell\_dfc.py}; geplanter Dualtest: \texttt{compute\_dual\_dfc.py}.
\end{remark}

\subsection{Erweiterungen}

\begin{itemize}
  \item \textbf{Kompressible Turbulenz:}
    Die Erweiterung auf kompressible Strömungen erfordert die Ersetzung
    der $L^2$-basierten Shell-Energien durch dichtegewichtete Grössen
    und die Modifikation der Closure zur Berücksichtigung akustischer
    Modenkopplung.
  \item \textbf{Magnetohydrodynamische (MHD) Turbulenz:}
    Das Iroshnikov--Kraichnan-Spektrum $E(k) \sim k^{-3/2}$ sollte aus
    einem modifizierten Funktional unter Einbeziehung von
    Alfv\'en-Wellen-Wechselwirkungen hervorgehen.
  \item \textbf{Zweidimensionale Turbulenz:}
    Die duale Kaskade (inverse Energiekaskade mit $k^{-5/3}$ und direkte
    Enstrophiekaskade mit $k^{-3}$~\cite{Kraichnan67}) stellt einen
    interessanten Testfall dar, der ein variationelles Problem mit zwei
    Nebenbedingungen erfordert.
\end{itemize}

\subsection{Post-Zookeeper-Split: Paper A schützen, Paper B schärfen}
\label{subsec:zookeeper-turbulence-transfer}

Der Zookeeper-Transfer klärt die Publikationsteilung.  Paper A sollte
die unbedingte Variationsaussage bleiben: K41 ist unter den angegebenen
Zulässigkeitsannahmen der eindeutige Minimierer des skalenaufgelösten
Freie-Energie-Funktionals.  Es sollte keine Claims über anomale
Dissipation, Onsager-Nicht-Eindeutigkeit oder physikalische
Lösungsselektion aufnehmen.

Auch hier ist der Post-Zookeeper-Import RFEP-v1.8-Normalformdisziplin,
kein neues unbedingtes Turbulenztheorem.  Der CORE-Abschluss betrifft den
Riemann-$C^{(2)}$/MS2-Schritt; Turbulenz muss seine eigenen Flux-,
Defekt-, Approximations-, Gap- und Viskositätsgrenzwerte beweisen.

Paper B ist der richtige Ort für die Defekttheorie.  Seine
Transferfelder sind:
\begin{description}[style=nextline,leftmargin=3.4cm,labelwidth=3.0cm]
  \item[Operator/Form.] Shell-Flux- und Littlewood--Paley-
  Energietransferoperatoren.
  \item[Defekt.] Der DFC-Flussdefekt, das anomalous-dissipation measure
  und die Abweichung von der K41-Steigung.
  \item[Approximation.] Shell-Modelle, DNS, Littlewood--Paley-Cutoffs und
  Viskositätsfamilien $\nu\to0$.
  \item[Gap.] Der Hessian-Gap des K41-Minimierers und der spektrale
  Poincare/SP-TU-Gap.
  \item[Grenzübergang.] Shell/DNS zur PDE, Viskosität gegen null und
  Langzeitmittel zur Inertial-Range-Aussage.
\end{description}

Das konkreteste Upgrade ist, DFC1 als duale gewichtete Bedingung
DFC1$^\vee$ über gleitende Skalenfenster zu formulieren. Dann kann
Kolmogorovs $4/5$-Gesetz zusammen mit der lokalen Energiebilanz von
Duchon--Robert und Eyink als exakter grobskaliger Energiefluss-Anker für
Paper B dienen, während die verbleibende Projektionsabschätzung auf
Littlewood--Paley-Gewichte explizit bleibt. Die Resultate von Brü--De
Lellis und Koautoren zur anomalen Dissipation liefern die Stresstest-Seite:
Der Defekt sollte als Maß oder Distribution formuliert werden, kompatibel
mit bekannten Vanishing-Viscosity-Szenarien.

\subsection{Offene Richtungen}

\begin{remark}[Energiekaskade und Minimax-Notation]
Das Kolmogorov-Spektrum lässt sich durch eine entkoppelte
Minimax-Penalty darstellen: Das Spektrum minimiert die KL-freie Energie,
während eine quadratische Penalty den Konstantfluss erzwingt. Das ist
nützliche Notation, aber kein bewiesenes dynamisches Nullsummenspiel
zwischen Energieeintrag und Dissipation.
\end{remark}

\begin{remark}[Dissipative Adaptation]
Englands Prinzip der dissipativen Adaptation legt nahe, dass getriebene Systeme sich selbst organisieren, um die Entropieproduktion zu maximieren. Im Turbulenzkontext liefert dies eine thermodynamische Begründung dafür, warum die Kaskade spontan das K41-Spektrum selektiert: Es ist dasjenige Profil, das die Rate der kinetischen Energiedissipation unter der Flussnebenbedingung maximiert.
\end{remark}

\begin{remark}[Kumulantenentwicklung für Intermittenz]
Die Positivität $\Phi \geq 0$ beschränkt die Kumulantenhierarchie der Energiefluktuationen: $|\kappa_n(E_j)| \leq n! \, \sigma^2 / (2\alpha^{n-2})$. Dies eröffnet einen direkten Zugang zu den Intermittenz-Exponenten $\zeta_p$ über die Hesse-Fluktuationen von $F[E]$ um den K41-Minimierer.
\end{remark}

% [Entfernt: Koalitionsbruch-Modell -- spekulative Metapher ohne formalen Inhalt]

% [Entfernt: Skalenabhängige-Kopplung-Analogie -- zu spekulativ für diesen Kontext]

\begin{remark}[Shell-Kaskade als Dobrushin-Einflussmatrix]
\label{rem:dobrushin-cascade}
Der schalenweise Energietransfer im Inertialbereich definiert eine
natürliche Dobrushin-Einflussmatrix $C = (c_{jk})$, wobei $c_{jk}$ den
Einfluss der Schale $k$ auf die Schale $j$ über den nichtlinearen
Transfer $T_j$ quantifiziert. Für die K41-Kaskade ergibt die
Nächste-Nachbar-Dominanz $c_{j,j\pm 1} = O(1)$ und
$c_{jk} = O(2^{-|j-k|})$ für $|j-k| > 1$. Der Spektralradius
$\rho(C) < 1$, also Dobrushin-Eindeutigkeit, ist äquivalent zur
Stabilität des K41-Minimierers: Kleine Störungen von $E^*$ an Schale
$k$ klingen exponentiell in $|j-k|$ ab.

Diese Perspektive, importiert aus dem Yang--Mills-Begleitpaper, in dem
die Dobrushin-Einflussmatrix die Gitter-Massenlücke kontrolliert,
offenbart einen strukturellen Isomorphismus: Die
Inertialbereichs-Kaskade ist ein ,,eindimensionales Gittersystem`` mit
exponentiell abklingenden Wechselwirkungen, und K41 ist sein eindeutiger
Gibbs-Zustand.
\end{remark}

\begin{remark}[DFC als defektive Monotonie mit kontrollierten Defekttermen]
\label{rem:dfc-defective-monotonicity}
Die Downhill-Flussbedingung lässt sich als \emph{defektive Monotonie}-Aussage verstehen, parallel zum defektiven Poincar\'{e}/LSI-Ansatz im Yang--Mills-Begleitpaper~\cite{Geiger2026-YM}:
\begin{itemize}
\item \emph{Idealfall (DFC1$^\vee$+DFC2+DFC3):} Das Freie-Energie-Funktional $F[E]$ ist streng monoton fallend entlang der Inertialkaskade. K41 ist der eindeutige Minimierer.
\item \emph{Defektiver Fall:} Intermittenz (Bottleneck-Effekt, lokale Rückstreuung) führt Defektterme $\delta_j$ an jeder Schale ein:
\[
F[E_{j+1}] \;\leq\; F[E_j] - \Delta_j + \delta_j, \qquad \sum_j \delta_j < \infty.
\]
Die Summierbarkeit der Defekte stellt sicher, dass die Kaskade trotz lokaler Verletzungen ,,asymptotisch abwärts'' bleibt. Dies ist strukturell identisch zum defektiven Dobrushin-Kriterium in Yang--Mills, wo der exponentiell unterdrückte Defekt $e^{-c\sqrt{\beta}}$ (Proposition~3.4 von~\cite{Geiger2026-YM}) dieselbe Rolle spielt wie die summierbaren $\delta_j$ hier.
\end{itemize}
Die Formulierung als defektive Monotonie macht die DFC \emph{numerisch testbar}: Man messe $\Delta_j$ und $\delta_j$ an jeder Schale in DNS-Daten und verifiziere, dass $\sum \delta_j / \sum \Delta_j \ll 1$. Die DFC-mit-Spielraum-Variante ist bis zu einem Datenfit- oder Perturbationstheorem am besten als Finite-Reynolds-Diagnostik zu lesen.
\end{remark}

\begin{remark}[Penalty-Sicht auf Inertialreduktionen]
\label{rem:minimax-inertial}
Wenn eine Inertialmannigfaltigkeits- oder Finite-Mode-Reduktion verfügbar
ist, kann die KL-Penalty in diesem reduzierten System mit expliziten
Closure-Fehlern getestet werden. Eine solche Reduktion könnte später einen
kontrollierten Rahmen für ein gekoppeltes Spiel liefern; dieses Paper
beweist jedoch nur die entkoppelte Penalty-Aussage aus
Proposition~\ref{prop:minimax}.
\end{remark}

% [Entfernt: Wirbelröhren als Quasiteilchen -- spekulativ, kein formales Ergebnis]

% [Entfernt: Abwärts im Knotenraum -- geometrische Metapher ohne rigorosen Inhalt]

% [Entfernt: Helizitäts-Kaskade-Brücke -- hängt von entfernter Knotenraum-Bemerkung ab]

\begin{remark}[Dobrushin-Diagnostik auf reduzierten Kaskaden]
\label{rem:dobrushin-inertial}
Die Shell-Kaskade als Dobrushin-Einflussmatrix
(Bemerkung~\ref{rem:dobrushin-cascade}) kann in Finite-Mode- oder
Inertialreduktionen getestet werden, sofern solche Reduktionen verfügbar
sind. Eine geprüfte Schranke $\rho(C)<1$ würde eine quantitative
Stabilitätsdiagnostik für die reduzierte Kaskade liefern, aber dieses
Paper beweist weder Dobrushin-Eindeutigkeit noch eine Turbulenz-Massenlücke.
Die Analogie zum Yang--Mills-Begleitpaper ist daher ein Forschungsprogramm,
kein importiertes Theorem.
\end{remark}

\begin{remark}[Pattern-A-Master-Stabilität -- problemübergreifende Struktur]
Diese Arbeit instanziiert das Pattern-A-Master-Stabilitätstemplate aus dem
vereinheitlichten Framework \cite{FST-Unified2026}: Strikte Konvexität der
renormierten freien Energie $F$ liefert eine lokale Stabilitätsmarge für
den K41-Minimierer. Eine echte Spektrallücke für einen Transferoperator
würde die oben offen gelassene zusätzliche Kaskadenoperator-Analyse
erfordern.
\end{remark}

% ===================================================================
\section*{Danksagung}

Der Autor dankt den KI-Systemen Claude und Gemini f\"ur Unterst\"utzung bei der Literaturrecherche, mathematischen Verifikation und Manuskripterstellung.

% ===================================================================
% ===================================================================
% KI-Offenlegung
% ===================================================================
\section*{KI-Offenlegung}
Bei der Erstellung dieses Manuskripts wurden KI-gestützte Werkzeuge (Claude, Anthropic)
in unterstützender Funktion eingesetzt, insbesondere für Literaturrecherche,
Textstrukturierung und Formulierungshilfe. Der konzeptionelle Entwurf,
die wissenschaftliche Argumentation und die Verantwortung für den gesamten Inhalt
liegen ausschließlich beim Autor.

\begin{thebibliography}{99}

\bibitem{K41a}
A.\,N.~Kolmogorov,
\emph{The local structure of turbulence in incompressible viscous fluid
for very large Reynolds numbers},
Doklady Akad.\ Nauk SSSR \textbf{30} (1941), 299--303.
Reprinted in Proc.\ R.\ Soc.\ Lond.\ A \textbf{434} (1991), 9--13.

\bibitem{K41b}
A.\,N.~Kolmogorov,
\emph{Dissipation of energy in the locally isotropic turbulence},
Doklady Akad.\ Nauk SSSR \textbf{32} (1941), 16--18.
Reprinted in Proc.\ R.\ Soc.\ Lond.\ A \textbf{434} (1991), 15--17.

\bibitem{K62}
A.\,N.~Kolmogorov,
\emph{A refinement of previous hypotheses concerning the local structure
of turbulence in a viscous incompressible fluid at high Reynolds number},
J.~Fluid Mech.\ \textbf{13} (1962), 82--85.

\bibitem{Onsager49}
L.~Onsager,
\emph{Statistical hydrodynamics},
Nuovo Cimento Suppl.\ \textbf{6} (1949), 279--287.

\bibitem{CET94}
P.~Constantin, W.~E, and E.\,S.~Titi,
\emph{Onsager's conjecture on the energy conservation for solutions of
Euler's equation},
Comm.\ Math.\ Phys.\ \textbf{165} (1994), 207--209.

\bibitem{DR00}
J.~Duchon and R.~Robert,
\emph{Inertial energy dissipation for weak solutions of incompressible
Euler and Navier--Stokes equations},
Nonlinearity \textbf{13} (2000), 249--255.

\bibitem{Isett18}
P.~Isett,
\emph{A proof of Onsager's conjecture},
Ann.\ of Math.\ \textbf{188} (2018), 871--963.

\bibitem{SL94}
Z.-S.~She and E.~L\'ev\^eque,
\emph{Universal scaling laws in fully developed turbulence},
Phys.\ Rev.\ Lett.\ \textbf{72} (1994), 336--339.

\bibitem{ES06}
G.\,L.~Eyink and K.\,R.~Sreenivasan,
\emph{Onsager and the theory of hydrodynamic turbulence},
Rev.\ Mod.\ Phys.\ \textbf{78} (2006), 87--135.

\bibitem{Frisch95}
U.~Frisch,
\emph{Turbulence: The Legacy of A.\,N.\ Kolmogorov},
Cambridge University Press, 1995.

\bibitem{Eyink03}
G.\,L.~Eyink,
\emph{Local $4/5$-law and energy dissipation anomaly in turbulence},
Nonlinearity \textbf{16} (2003), 137--145.

\bibitem{Kraichnan67}
R.\,H.~Kraichnan,
\emph{Inertial ranges in two-dimensional turbulence},
Phys.\ Fluids \textbf{10} (1967), 1417--1423.

\bibitem{MY75}
A.\,S.~Monin and A.\,M.~Yaglom,
\emph{Statistical Fluid Mechanics}, Bde.~1--2,
MIT Press, 1971--1975.

\bibitem{Kaneda03}
Y.~Kaneda, T.~Ishihara, M.~Yokokawa, K.~Itakura, and A.~Uno,
\emph{Energy dissipation rate and energy spectrum in high resolution
direct numerical simulations of turbulence in a periodic box},
Phys.\ Fluids \textbf{15} (2003), L21--L24.

\bibitem{Gotoh02}
T.~Gotoh, D.~Fukayama, and T.~Nakano,
\emph{Velocity field statistics in homogeneous steady turbulence
obtained using a high-resolution direct numerical simulation},
Phys.\ Fluids \textbf{14} (2002), 1065--1081.

\bibitem{Pope00}
S.\,B.~Pope,
\emph{Turbulent Flows},
Cambridge University Press, 2000.

\bibitem{FST-Unified2026}
L.~Geiger, \emph{The Renormalized Free Energy Framework: A Unified Resolution of Structural Bottlenecks},
Zenodo, March 2026.
\href{https://doi.org/10.5281/zenodo.19036190}{doi:10.5281/zenodo.19036190}.

\bibitem{BonnansShapiro00}
J.\,F.~Bonnans und A.~Shapiro,
\emph{Perturbation Analysis of Optimization Problems},
Springer Series in Operations Research, Springer, New York, 2000.

\bibitem{Rockafellar70}
R.\,T.~Rockafellar,
\emph{Convex Analysis},
Princeton Mathematical Series, Bd.~28,
Princeton University Press, 1970.

\bibitem{FoiasSellTemam88}
C.~Foias, G.\,R.~Sell und R.~Temam,
\emph{Inertial manifolds for nonlinear evolutionary equations},
J.~Differential Equations \textbf{73} (1988), 309--353.

\bibitem{Moffatt69}
H.\,K.~Moffatt,
\emph{The degree of knottedness of tangled vortex lines},
J.~Fluid Mech.\ \textbf{35} (1969), 117--129.

\bibitem{ArnoldKhesin98}
V.\,I.~Arnold und B.\,A.~Khesin,
\emph{Topological Methods in Hydrodynamics},
Applied Mathematical Sciences, Bd.~125,
Springer, New York, 1998.

\bibitem{Drivas2019cascade}
T.\,D.~Drivas,
\emph{Turbulent cascade direction and Lagrangian time-asymmetry},
J.~Nonlinear Sci.\ \textbf{29} (2019), 65--88.
\href{https://arxiv.org/abs/1802.02289}{arXiv:1802.02289}.

\bibitem{DeRosaDrivasInversi2024}
L.~De~Rosa, T.\,D.~Drivas und M.~Inversi,
\emph{On the support of anomalous dissipation measures},
J.~Math.\ Fluid Mech.\ \textbf{26} (2024), Art.~62.
\href{https://arxiv.org/abs/2301.09603}{arXiv:2301.09603}.

\bibitem{AlbrittonBrueColombo2022}
D.~Albritton, E.~Bru\'e und M.~Colombo,
\emph{Non-uniqueness of Leray solutions of the forced Navier--Stokes equations},
Ann.\ of Math.\ \textbf{196} (2022), 415--455.

\bibitem{Geiger2026-YM}
L.~Geiger, \textit{Volume-Independent Spectral Gap for Strong-Coupling Lattice Gauge Theory via Poincar\'e-Dobrushin Machinery}, Preprint, 2026.

\bibitem{SaddoughiVeeravalli94}
S.\,G.~Saddoughi und S.\,V.~Veeravalli,
\emph{Local isotropy in turbulent boundary layers at high Reynolds number},
J.~Fluid Mech.\ \textbf{268} (1994), 333--372.

\bibitem{JHTDB08}
Y.~Li, E.~Perlman, M.~Wan, Y.~Yang, C.~Meneveau, R.~Burns, S.~Chen,
A.~Szalay und G.~Eyink,
\emph{A public turbulence database cluster and applications to study
Lagrangian evolution of velocity increments in turbulence},
J.~Turbulence \textbf{9} (2008), Nr.~31.

\bibitem{Lvov98}
V.\,S.~L'vov, E.~Podivilov, A.~Pomyalov, I.~Procaccia
und D.~Vandembroucq,
\emph{Improved shell model of turbulence},
Phys.\ Rev.\ E \textbf{58} (1998), 1811--1822.

\bibitem{BV19}
T.~Buckmaster und V.~Vicol,
\emph{Nonuniqueness of weak solutions to the Navier--Stokes equation},
Ann.\ of Math.\ \textbf{189} (2019), 101--144.

\bibitem{England2015}
J.\,L.~England,
\emph{Dissipative adaptation in driven self-assembly},
Nature Nanotechnology \textbf{10} (2015), 919--923.

\end{thebibliography}

\end{document}
